%%%%%%%%%%%%%%%%%%%%%%%%%%%%%%%%%%%%%%%%%%%%%%%%%%%%%%%%%%%%%%%%%%%%%%%%%%%%%%
%  Marketing Management (IM345AT) -- Volume 3 of 5
%  UNIT III : Web Analytics, Google Analytics and E-mail Marketing
%%%%%%%%%%%%%%%%%%%%%%%%%%%%%%%%%%%%%%%%%%%%%%%%%%%%%%%%%%%%%%%%%%%%%%%%%%%%%%

\documentclass{MMTextbook}

\unitnumber{III}
\unittitle{Web Analytics, Google Analytics and E-mail Marketing}
\unitsubtitle{Measurement, reporting, and the owned channel with the highest
conversion rate in digital marketing}
\unithours{07 Hours}
\volumeno{3}

\makeindex[intoc=true,columns=2,title=Index]

\hypersetup{
  pdftitle={Marketing Management IM345AT -- Unit III: Web Analytics, Google
             Analytics and E-mail Marketing},
  pdfsubject={Semester IV, Department of Industrial Engineering and Management,
              RV College of Engineering},
  pdfkeywords={web analytics, Google Analytics, e-mail marketing, conversion,
               IM345AT, marketing management}
}

\begin{document}

%%%%%%%%%%%%%%%%%%%%%%%%%%%%%%%%%%%%%%%%%%%%%%%%%%%%%%%%%%%%%%%%%%%%%%%%%%%%%%
\frontmatter
\pagestyle{mmfront}

\halftitlepage
\maintitlepage

\colophonpage{%
\textbf{\coursetitle\ ---\ The Study Text}\\[2pt]
Volume 3 of 5 \textbullet\ Unit III: Web Analytics, Google Analytics and E-mail
Marketing\\[10pt]
%
Prepared for the course \coursetitle\ (\coursecode), Semester IV, of the
\coursedept, \courseinstitute.\\[10pt]
%
\textbf{Sources.} The syllabus, course outcomes, assessment rubrics and
reference-book list reproduced in the front matter are transcribed from the
official course document. The previous-year questions in Appendix~B are
transcribed from RVCE Continuous Internal Evaluation and Semester End
Examination papers for this course, and where an official scheme of valuation
was released the model answers follow it. The explanatory text is compiled from
classroom lecture notes and, where those were incomplete, from the prescribed
reference texts.\\[10pt]
%
\textbf{Typography.} Set in Nimbus Roman with URW Gothic display headings.
Composed in \LaTeX{} using the \texttt{MMTextbook} class written for this
series. The visual design is derived from \emph{The Legrand Orange Book}
template by Mathias Legrand, with modifications by Vel, distributed through
LaTeXTemplates.com under a Creative Commons BY-NC-SA 4.0 licence; the class
re-implements that design language on standard CTAN packages. All figures are
drawn in \TikZ{} and \texttt{pgfplots}.\\[10pt]
%
\textbf{Status.} Unofficial student-compiled study material. Not issued,
endorsed or verified by the Department or by RV College of Engineering.\\[10pt]
%
2026 edition.}

% ---------------------------------------------------------------- preface
\frontchapter{Preface to Volume 3}

Unit~III is where digital marketing stops being persuasion and becomes
measurement. The two halves of the unit are the instrument and the channel: web
analytics is how a business finds out what actually happened, and e-mail
marketing is the one channel whose audience the business owns outright.

The connection between them is not incidental. Web analytics converts anonymous
traffic into evidence, and without it digital marketing forfeits the single
structural advantage it holds over traditional marketing --- measurability. E-mail
is the channel that most rewards that evidence, because it is the only major
channel where the marketer can segment, personalise, trigger and re-send without
asking an algorithm's permission.

Chapters~1 to~3 cover web analytics and Google Analytics; Chapters~4 to~7 cover
e-mail marketing. The chapter sequence follows the syllabus phrase order exactly,
and Appendix~A records the mapping.

One characteristic of this unit deserves stating in advance, because it changes
how the material should be learned. Almost every past-paper question in Unit~III
is \emph{applied}: you are given a business --- an e-commerce site, an e-learning
platform, a travel agency, a skincare startup, an online business targeting
college students --- and asked to build a plan or diagnose a failure. General
commentary scores badly on such questions. What scores is a named framework
applied to the specific case. This volume therefore states each framework
explicitly and by name --- the six-stage measurement funnel, the seven-node
e-mail plan, the eight-step goal setup, the eight-step segmentation guide --- and
then works it through on a real scenario so the pattern is visible.

Unlike Volumes~1 and~2, this volume has no appendix of off-syllabus material.
Every topic in the source notes for Unit~III maps to a phrase in the unit
descriptor, so nothing had to be set aside.

\vspace{12pt}
\noindent\emph{The whole of the course syllabus is reproduced overleaf. The
unit covered by this volume is Unit~III.}

% ------------------------------------------------- conventions & syllabus
%%%%%%%%%%%%%%%%%%%%%%%%%%%%%%%%%%%%%%%%%%%%%%%%%%%%%%%%%%%%%%%%%%%%%%%%%%%%%%
%  mm-howtouse.tex
%  Editorial conventions for the series.  Shared, identical, across all five
%  volumes so that a reader who learns the apparatus once can use any volume.
%%%%%%%%%%%%%%%%%%%%%%%%%%%%%%%%%%%%%%%%%%%%%%%%%%%%%%%%%%%%%%%%%%%%%%%%%%%%%%

\frontchapter{How to Use This Study Text}

\section*{The shape of a volume}
\addcontentsline{toc}{section}{The shape of a volume}

Each of the five volumes in this series covers exactly one unit of the
IM345AT syllabus and is self-contained. Every volume opens with the official
syllabus for the whole course, so that the unit in hand can always be located
within it, and closes with a set of appendices carrying the concordance,
the question bank, the glossary and a quick-reference sheet.

The numbered chapters of each volume map one-to-one onto the topic phrases of
that unit's official syllabus descriptor. The order of the chapters is the
order of the syllabus. Appendix~A sets out that mapping explicitly, phrase by
phrase, so coverage can be audited rather than assumed.

\section*{The five panel types}
\addcontentsline{toc}{section}{The five panel types}

Material that is not continuous prose is set in one of five panels. Each has a
fixed meaning throughout the series.

\begin{definitionbox}[the term being defined]
A formal definition, worth reproducing close to verbatim in an answer script.
The term appears in the panel heading and again in the index.
\end{definitionbox}

\begin{keybox}[High Yield]
A result, list or distinction that carries disproportionate examination weight.
Where a topic has recurred across papers, the recurrence is recorded as a
bracketed tag. \pyq{CIE-I, 3 Apr 2025; CIE-I, 16 Apr 2026}
\end{keybox}

\begin{exambox}[Model Answer]
An answer written at the length and in the register the marking scheme
rewards. Where an official scheme of valuation was released, the panel follows
it.
\end{exambox}

\begin{examplebox}[Worked Example]
A concrete illustration --- a named company, a named persona, a worked
calculation --- rather than a general statement.
\end{examplebox}

\begin{warnbox}
A common and costly error: a distinction usually collapsed, a term usually
misused, or an answer that looks complete and is not.
\end{warnbox}

A sixth panel, in green, sets out a framework or process as a single
sequence:

\begin{frameworkbox}[Framework]
Step one \ra\ step two \ra\ step three \ra\ back to step one.
\end{frameworkbox}

\section*{How previous-year questions are recorded}
\addcontentsline{toc}{section}{How previous-year questions are recorded}

\begin{keybox}[The Bracket Convention]
A past-paper appearance is recorded \textbf{only} as a bracketed source tag
attached to the material, never as a sentence of narrative. So the text reads
\begin{quote}\sffamily\footnotesize
Push versus pull marketing. \pyq{CIE-I, 3 Apr 2025; CIE-I, 16 Apr 2026; SEE, Jun/Jul 2025 --- 2 M}
\end{quote}
and never ``this was asked in 2025 and again in 2026''. The tag carries the
paper, the date and, where the mark allocation is known, the marks. Multiple
appearances are separated by semicolons.
\end{keybox}

\noindent The abbreviations used inside tags are:

\vspace{4pt}
{\sffamily\small
\begin{tabular}{@{}P{0.14\textwidth}P{0.80\textwidth}@{}}
\toprule[1.1pt]
\rlab{CIE-I} & First Continuous Internal Evaluation test\\
\rlab{CIE-II} & Second Continuous Internal Evaluation test\\
\rlab{CIE-III} & Third Continuous Internal Evaluation test\\
\rlab{SEE} & Semester End Examination\\
\rlab{2 M} & Mark allocation of the question as printed on the paper\\
\bottomrule[1.1pt]
\end{tabular}}

\vspace{10pt}

Appendix~B of each volume reproduces every past-paper question traced to that
unit, in full, grouped by mark weight, each carrying the same bracketed tag.

\section*{Cross-references, figures and the index}
\addcontentsline{toc}{section}{Cross-references, figures and the index}

Cross-references are given as \S\,chapter.section --- so \S\,1.4 is the fourth
section of Chapter~1. Every figure in this series is drawn in \TeX{} rather
than imported as an image, so the diagrams print at full resolution at any
size and use the same typeface as the surrounding text. Figures and tables are
listed after the table of contents.

Terms set in \textbf{\textit{bold italic}} on first use are indexed. The index
at the back of each volume is a real back-of-book index built from those
terms, not a repeat of the contents list.

\section*{Status of this text}
\addcontentsline{toc}{section}{Status of this text}

\begin{warnbox}[Unofficial Document]
This is a student-compiled study text. It is not issued, endorsed or verified
by the Department of Industrial Engineering and Management or by RV College of
Engineering. The syllabus reproduced in the front matter, the assessment
rubrics and the previous-year questions are transcribed from official
documents; the explanatory prose around them is not official. Where the source
lecture notes were incomplete, gaps have been filled from the prescribed
reference texts listed in the front matter. Always cross-check against the
official syllabus and your own class notes.
\end{warnbox}

%%%%%%%%%%%%%%%%%%%%%%%%%%%%%%%%%%%%%%%%%%%%%%%%%%%%%%%%%%%%%%%%%%%%%%%%%%%%%%
%  mm-syllabus.tex
%  Verbatim reproduction of the official IM345AT course document.
%  Shared, byte-identical, across all five volumes of the series.
%%%%%%%%%%%%%%%%%%%%%%%%%%%%%%%%%%%%%%%%%%%%%%%%%%%%%%%%%%%%%%%%%%%%%%%%%%%%%%

\frontchapter{Official Syllabus: IM345AT}

\noindent
{\sffamily\small\color{slatelight}%
Reproduced from the official course document issued by the Department of
Industrial Engineering and Management, RV College of Engineering. The text of
the five unit descriptors below is the authority against which every chapter of
this series is written; nothing outside it appears in the numbered chapters of
any volume.\par}

\vspace{10pt}

\begin{center}
\sffamily\small
\begin{tabular}{@{}P{0.238\textwidth}P{0.238\textwidth}P{0.218\textwidth}P{0.218\textwidth}@{}}
\toprule[1.1pt]
\multicolumn{4}{@{}l}{\bfseries\normalsize Semester: IV \quad\textbullet\quad MARKETING MANAGEMENT}\\
\multicolumn{4}{@{}l}{\color{slatelight}Category: Professional Core Course (Theory)}\\
\midrule
\rlab{Course Code} & IM345AT & \rlab{CIE Marks} & 100 Marks\\
\rlab{Credits (L:T:P)} & 3:0:0 & \rlab{SEE Marks} & 100 Marks\\
\rlab{Total Hours} & 45 L & \rlab{SEE Duration} & 3.00 Hours\\
\bottomrule[1.1pt]
\end{tabular}
\end{center}

\vspace{6pt}

% ---------------------------------------------------------------- unit blocks
\newcommand{\syllunit}[3]{%
  \begin{tcolorbox}[mmbase, colback=white, colframe=ocre!45, boxrule=0.7pt,
      sharp corners, fontupper=\small,
      fonttitle=\sffamily\bfseries\small\color{white},
      coltitle=white, colbacktitle=#3,
      title={#1},
      attach boxed title to top left={xshift=-0.7pt, yshift=-0.7pt},
      boxed title style={sharp corners, boxrule=0pt, colback=#3},
      top=16pt]
    #2
  \end{tcolorbox}}

\syllunit{UNIT--I \hfill 07 Hrs}{%
\textbf{Introduction to Digital Marketing:} Principles of Digital Marketing;
Digital Marketing Channels; Tools to Create Buyer Persona; Competitor Research
Tools, Website Analysis Tools, etc.\par
\vspace{3pt}
\textbf{Content Marketing:} Content Marketing Concepts \& Strategies; Planning,
Creating, Distributing \& Promoting Content; Optimize Website UX \& Landing
Pages; Measure Impact; Metrics \& Performance; Using Content Research for
Opportunities, etc.}{ocre}

\syllunit{UNIT--II \hfill 08 Hrs}{%
\textbf{Social Media Marketing:} Introduction; Major Social Media Platforms for
Marketing; Developing Data-driven Audience \& Campaign Insights; Social Media
for Business; Creation \& Optimization of Social Media Campaigns, etc.\par
\vspace{3pt}
\textbf{Search Engine Optimization:} Search Engine Optimization Fundamentals;
Keywords and SEO Content Plan; SEO \& Business Objectives; Writing SEO Content;
On-site \& off-site SEO; Optimize Organic Search Ranking, etc.}{ocre}

\syllunit{UNIT--III \hfill 07 Hrs}{%
\textbf{Web Analytics \& Google Analytics:} Google Analytics Tools; Web
Analytics Tools, etc.\par
\vspace{3pt}
\textbf{E-mail Marketing:} Effective E-mail Campaigns; E-mail Plan; E-mail
Marketing Campaign Analysis; Measuring Conversions \& keeping up, etc.}{ocre}

\syllunit{UNIT--IV \hfill 07 Hrs}{%
\textbf{Web Design:} Web design, optimization of websites; Publishing a basic
website; User-centered Design and Website Optimization; Design Principles and
Website Copy; Website Metrics \& Developing Insight, etc.\par
\vspace{3pt}
\textbf{Mobile Marketing:} Difference between mobile advertising and marketing,
utilizing mobile marketing for sales promotions, online applications,
etc.}{ocre}

\syllunit{UNIT--V \hfill 07 Hrs}{%
\textbf{Conversion Optimization:} What is AIDAS and its role; website
optimization; what visitors want to see on the website; how to optimize key
element and increase the effect of landing on a particular page\par
\vspace{3pt}
\textbf{Digital Analytics:} Evolution of Digital Analytics, information about
end-to-end customer experience, analyst's influence on business, role as a
change agent, etc.}{ocre}

% ------------------------------------------------------------ course outcomes
\section*{Course Outcomes}
\addcontentsline{toc}{section}{Course Outcomes}

\noindent{\small After completing the course, the students will be able to:\par}
\vspace{4pt}

{\sffamily\small
\begin{longtable}{@{}P{0.07\textwidth}P{0.89\textwidth}@{}}
\toprule[1.1pt]
\rlab{CO1} & Differentiate the benefits drawn by updated marketing mix from
traditional marketing mix for effective marketing management there by to stay
competitive in today's global market-place.\\
\midrule
\rlab{CO2} & Develop an effective holistic marketing atmosphere to efficiently
face the challenges in dynamically changing market.\\
\midrule
\rlab{CO3} & Formulate a potential marketing plan to effectively reach the
targeted market segments, by delivering the value to targeted customers through
practicing sound marketing research.\\
\midrule
\rlab{CO4} & Create new channels to improvise marketing to achieve and maintain
competitive position in globalized market-place.\\
\bottomrule[1.1pt]
\end{longtable}}

% --------------------------------------------------------------- reference books
\section*{Prescribed Reference Books}
\addcontentsline{toc}{section}{Prescribed Reference Books}

{\small
\begin{enumerate}[leftmargin=2.1em]
\item \textit{Marketing Management}, Philip Kotler, Kevin Lane Keller, 15th
      Edition, 2016, Pearson, ISBN: 978-93-325-5718-5.
\item \textit{Digital Marketing --- Strategy, Implementation \& Practice}, Dave
      Chaffey, Fiona Ellis-Chadwick, 7th Edition, 2019, Pearson,
      ISBN: 9781292241623, 1292241624.
\item \textit{Marketing Research}, Donald S. Tull, Del I. Hawkins, 6th Edition,
      Prentice Hall India, ISBN: 8120309618.
\item \textit{Marketing Management --- A South Asian Perspective}, Philip
      Kotler, Kevin Lane Keller, Abraham Koshy, Mithileshwar Jha, 14th Edition,
      2013, Pearson, ISBN: 978-81-317-6716-0.
\item \textit{Marketing Research}, David A. Aaker, V. Kumar, George S. Day, 9th
      Edition, 2008, John Wiley \& Sons, ISBN: 978-265-1791-6.
\end{enumerate}}

% ------------------------------------------------------------------- rubrics
\frontchapter{Assessment Scheme}

\section*{Rubric for the Continuous Internal Evaluation (Theory)}
\addcontentsline{toc}{section}{Rubric for the Continuous Internal Evaluation}

{\sffamily\small
\begin{longtable}{@{}P{0.045\textwidth}P{0.80\textwidth}>{\raggedleft\arraybackslash}P{0.07\textwidth}@{}}
\toprule[1.1pt]
\rowcolor{ocre}\thd{\#} & \thd{Component} & \thd{Marks}\\
\midrule
\rlab{1.} & \textbf{Quizzes:} Quizzes will be conducted in online/offline mode.
TWO QUIZZES will be conducted \& each quiz will be evaluated for 10 marks. The
sum of two quizzes will be the final quiz marks. & \bfseries 20\\
\midrule
\rowcolor{tablealt}
\rlab{2.} & \textbf{Tests:} Students will be evaluated in test, descriptive
questions with different complexity levels (Revised Bloom's Taxonomy Levels:
Remembering, Understanding, Applying, Analyzing, Evaluating and Creating). TWO
tests will be conducted. Each test will be evaluated for 50 marks, adding up to
100 marks. Final test marks will be reduced to 40 marks. & \bfseries 40\\
\midrule
\rlab{3.} & \textbf{Experiential Learning:} Students will be evaluated for their
creativity and practical implementation of the problem. Case study-based
teaching learning (10), Program specific requirements (10), Video based
seminar/presentation/demonstration (20), adding up to 40. & \bfseries 40\\
\midrule
\multicolumn{2}{@{}l}{\bfseries MAXIMUM MARKS FOR THE CIE THEORY} & \bfseries 100\\
\bottomrule[1.1pt]
\end{longtable}}

\section*{Rubric for the Semester End Examination (Theory)}
\addcontentsline{toc}{section}{Rubric for the Semester End Examination}

{\sffamily\small
\begin{longtable}{@{}P{0.12\textwidth}P{0.72\textwidth}>{\raggedleft\arraybackslash}P{0.08\textwidth}@{}}
\toprule[1.1pt]
\rowcolor{ocre}\thd{Q. No.} & \thd{Contents} & \thd{Marks}\\
\midrule
\multicolumn{3}{@{}l}{\bfseries\color{slate} PART A}\\
\midrule
\rlab{1} & Objective type questions covering entire syllabus & \bfseries 20\\
\midrule
\multicolumn{3}{@{}l}{\bfseries\color{slate} PART B \quad\normalfont\itshape(Maximum of TWO sub-divisions only)}\\
\midrule
\rlab{2}      & Unit 1: (Compulsory)      & \bfseries 16\\
\rowcolor{tablealt}
\rlab{3 \& 4} & Unit 2: (Internal Choice) & \bfseries 16\\
\rlab{5 \& 6} & Unit 3: (Internal Choice) & \bfseries 16\\
\rowcolor{tablealt}
\rlab{7 \& 8} & Unit 4: (Internal Choice) & \bfseries 16\\
\rlab{9 \& 10}& Unit 5: (Internal Choice) & \bfseries 16\\
\midrule
\multicolumn{2}{@{}l}{\bfseries TOTAL} & \bfseries 100\\
\bottomrule[1.1pt]
\end{longtable}}

\vspace{8pt}

\begin{keybox}[What the SEE Structure Means for This Volume]
Part A draws twenty marks of objective questions from the \emph{entire}
syllabus, so no unit can be skipped. In Part B, Question~2 is compulsory and is
always set from Unit~I; Units~II to~V each carry a 16-mark internal-choice
pair. A candidate therefore answers Unit~I plus one question from each of the
remaining four units --- five full questions in all.
\end{keybox}


% ---------------------------------------------------------------- contents
\cleardoublepage
\pdfbookmark[0]{Contents}{toc}
\tableofcontents
\cleardoublepage
\listoffigures
\cleardoublepage
\listoftables

%%%%%%%%%%%%%%%%%%%%%%%%%%%%%%%%%%%%%%%%%%%%%%%%%%%%%%%%%%%%%%%%%%%%%%%%%%%%%%
\mainmatter
\pagestyle{mmmain}

%%%%%%%%%%%%%%%%%%%%%%%%%%%%%%%%%%%%%%%%%%%%%%%%%%%%%%%%%%%%%%%%%%%%%%%%%%%%%%
\chapter{Web Analytics --- Fundamentals and Role}
%%%%%%%%%%%%%%%%%%%%%%%%%%%%%%%%%%%%%%%%%%%%%%%%%%%%%%%%%%%%%%%%%%%%%%%%%%%%%%

\syllabusstrip{Web Analytics \& Google Analytics.}

\begin{learningoutcomes}
\begin{itemize}
\item What web analytics is, and why it matters in digital marketing.
\item Traditional marketing analysis against web analytics-based analysis.
\item The three metric families: volume, quality, action.
\item The core metrics, precisely defined --- including the distinction between
      bounce\ix{bounce} rate and exit rate, and between dimensions and metrics.
\item The dual lens of website analytics and social analytics.
\item The eight roles of web analytics and the seven benefits of analysing site
      data.
\item How to align analytics KPIs with organisational goals.
\item A/B testing\ix{A/B testing} and conversion rate optimisation.
\end{itemize}
\end{learningoutcomes}

\section{Definition and Importance}

\begin{definitionbox}[Web Analytics]
\keytermas{Web analytics}{web analytics} is the process of collecting,
measuring, analysing and reporting website data in order to understand and
optimise web usage.
\pyq{CIE-II, 26 May 2026 --- 2 M}
\end{definitionbox}

\noindent\textbf{Importance in digital marketing:} it helps understand customer
behaviour; measures campaign effectiveness; improves website performance; and
supports better marketing decisions.

\begin{keybox}[Why This Is the Foundation of the Whole Discipline]
Web analytics converts anonymous traffic into evidence. Without it, digital
marketing loses the one structural advantage it holds over traditional
marketing --- measurability.
\end{keybox}

\section{Traditional Marketing Analysis Against Web Analytics}

\begin{keybox}[High Yield]
\pyq{CIE-II, 26 May 2026 --- 10 M}
Learn at least seven rows, and attach an example to each.
\end{keybox}

\begin{mmlongtable}{Traditional marketing analysis against web analytics-based analysis}{tab:tradvsweb}%
{@{}P{0.16\textwidth}P{0.36\textwidth}P{0.41\textwidth}@{}}
\rowcolor{ocre}\thd{Aspect} & \thd{Traditional Marketing Analysis} & \thd{Web Analytics-Based Analysis}\\
\midrule
\endfirsthead
\rowcolor{ocre}\thd{Aspect} & \thd{Traditional Marketing Analysis} & \thd{Web Analytics-Based Analysis}\\
\midrule
\endhead
\rlab{Data collection} & Surveys and interviews & Real-time online tracking\\
\rowcolor{tablealt}
\rlab{Speed} & Slow & Instant\\
\rlab{Accuracy} & Limited --- relies on recall and self-reporting & Highly
accurate --- records actual behaviour\\
\rowcolor{tablealt}
\rlab{Customer tracking} & Difficult & Easy\\
\rlab{Cost} & Expensive & Cost-effective\\
\rowcolor{tablealt}
\rlab{Measurement} & Estimated results & Measurable KPIs\\
\rlab{Sample} & A sample of respondents & The entire population of visitors\\
\rowcolor{tablealt}
\rlab{Attribution} & Largely guesswork & Channel- and campaign-level
attribution\\
\rlab{Feedback loop} & Post-campaign & Continuous, mid-campaign\\
\rowcolor{tablealt}
\rlab{Example} & A focus group is asked whether an advertisement was memorable &
Google Analytics shows exactly how many people saw the advertisement, clicked,
browsed and purchased\\
\end{mmlongtable}

\section{The Three Metric Families}

\begin{figure}[htbp]
\centering
\begin{tikzpicture}[font=\sffamily\footnotesize,
    fam/.style={rounded corners=4pt,line width=1pt,align=center,
                text width=3.35cm,inner sep=8pt,font=\sffamily\scriptsize}]
  \node[fam,draw=ocre,fill=ocrepale,text=slate] (v) at (0,0)
    {\textbf{\textcolor{ocredark}{VOLUME}}\\[3pt]
     \textit{Are we being seen?}\ \ (attract)\\[5pt]
     impressions\\ visits and sessions\\ reach};
  \node[fam,draw=deepblue,fill=deepbluepale,text=slate] (q) at (4.05,0)
    {\textbf{\textcolor{deepblue}{QUALITY}}\\[3pt]
     \textit{Are we relevant?}\ \ (engage)\\[5pt]
     bounce rate\\ session duration\\ sentiment};
  \node[fam,draw=greenok,fill=greenpale,text=slate] (a) at (8.1,0)
    {\textbf{\textcolor{greenok}{ACTION}}\\[3pt]
     \textit{Are they engaging?}\ \ (measure)\\[5pt]
     click-through rate\\ conversion rate\\ engagement rate};
  \draw[-{Latex[length=2.4mm]},slate,line width=0.9pt] (v)--(q);
  \draw[-{Latex[length=2.4mm]},slate,line width=0.9pt] (q)--(a);
\end{tikzpicture}
\caption[The three metric families]%
        {The three families in sequence. A reading in one family cannot be
         interpreted without the others: high volume with poor quality is wasted
         acquisition, and good quality with no action is a missing call to
         action.}
\label{fig:metricfamilies}
\end{figure}

\begin{mmlongtable}{The three metric families}{tab:metricfam3}%
{@{}P{0.13\textwidth}P{0.24\textwidth}P{0.56\textwidth}@{}}
\rowcolor{ocre}\thd{Family} & \thd{Question it answers} & \thd{Metrics}\\
\midrule
\endfirsthead
\rowcolor{ocre}\thd{Family} & \thd{Question it answers} & \thd{Metrics}\\
\midrule
\endhead
\rlab{Volume} & Are we being seen? (attract) & \textbf{Impressions} --- the
number of times content is seen. \textbf{Visits or sessions} --- total times
people have been to the site. \textbf{Reach} --- total users receiving brand
content.\\
\rowcolor{tablealt}
\rlab{Quality} & Are we relevant? (engage) & \textbf{Bounce rate} --- people
arriving and leaving without visiting another page. \textbf{Session duration}
--- time spent per visit. \textbf{Sentiment} --- positive, negative or neutral
brand mentions.\\
\rlab{Action} & Are they engaging? (measure) & \textbf{Click-through rate} ---
the percentage of people who see content and click. \textbf{Conversion rate}
--- the percentage who buy against who arrived. \textbf{Engagement rate} ---
the ratio of people taking action to the total who saw the content.\\
\end{mmlongtable}

\section{Core Metrics, Precisely Defined}

\begin{description}
\item[Users or unique visitors] The number of distinct individuals who visited
the site in a period, de-duplicated by client identifier.
\item[Sessions or visits] A group of interactions by one user within a given time
frame. A single user may generate many sessions.
\item[Pageviews] The total number of pages loaded. \emph{Unique} pageviews counts
a page once per session.
\item[Bounce rate] The percentage of visitors who leave a website after viewing
only one page without any interaction.
\item[Exit rate] The percentage of sessions that ended on a given page.
\item[Average session duration] The mean time users spend per session; a proxy
for content depth and interest.
\item[Pages per session] The average number of pages viewed per session; a
measure of site exploration.
\item[Conversion rate] The percentage of sessions or users completing a defined
goal --- purchase, sign-up, download, enquiry.
\item[Goal or conversion] A specific, valuable, trackable user action that the
site is designed to produce.
\item[Acquisition channel] The grouping that describes where a session came from
--- organic search\ix{organic search}, paid search, direct, referral, social, e-mail, display.
\item[New against returning visitors] The split between first-time and repeat
users; a retention signal.
\end{description}

\begin{warnbox}[Bounce Rate Is Not Exit Rate]
Every session that bounces also exits, but not every exit is a bounce. A bounce
is a single-page session with no interaction; an exit is simply the last page of
\emph{any} session.

\smallskip
The practical consequence: an exit is not necessarily a failure. An
order-confirmation page \emph{should} have a high exit rate, because the user has
finished. Treating exit rate as a fault indicator therefore leads to
`optimising' pages that are already working.
\end{warnbox}

\begin{definitionbox}[Dimensions and Metrics]
\keytermas{Dimensions}{dimension} are the \emph{attributes} of the data --- the
qualitative descriptors that answer ``what'' or ``where'': city, device, browser,
source and medium, page path, landing page.

\keytermas{Metrics}{metric} are the \emph{quantitative} measurements --- the
numbers: sessions, users, pageviews, bounce rate, revenue.

\smallskip
In a report, dimensions form the rows and metrics form the columns.
\pyq{SEE, Oct/Nov 2023 --- 2 M}
\end{definitionbox}

\section{Bounce Rate and Website Performance Assessment}

\begin{exambox}[Model Answer --- 2 Marks]
\pyq{CIE-II, 26 May 2026; SEE, Jun/Jul 2025 --- 2 M}
Bounce rate is the percentage of visitors who leave a website after viewing only
one page without any interaction.

\smallskip
\textbf{Effect on performance evaluation.} A high bounce rate indicates that
visitors are not finding what they expected --- signalling poor content
relevance, slow page load, confusing design, weak calls to action or mismatched
traffic. It therefore acts as an early diagnostic of both \emph{content quality}
and \emph{traffic quality}, and correlates with lower conversion rates.
\end{exambox}

\begin{warnbox}[Always Read Bounce Rate in Context]
A high bounce is not automatically a fault. A blog post that answers a question
completely, or a contact page that supplies a phone number, may legitimately show
a high bounce rate because the visitor's need was met on that single page. The
diagnosis depends on what the page was \emph{for}.
\end{warnbox}

\section{The Dual Lens --- Website Analytics and Social Analytics}

\begin{mmlongtable}{The website footprint and the social echo}{tab:duallens}%
{@{}P{0.13\textwidth}P{0.42\textwidth}P{0.38\textwidth}@{}}
\rowcolor{ocre}\thd{} & \thd{The Website Footprint} & \thd{The Social Echo}\\
\midrule
\endfirsthead
\rowcolor{ocre}\thd{} & \thd{The Website Footprint} & \thd{The Social Echo}\\
\midrule
\endhead
\rlab{Focus} & On-site behaviour and technical specifications & Off-site
popularity and brand conversation\\
\rowcolor{tablealt}
\rlab{Key data} & Device type, make and model; browser, resolution, operating
system; language and location; new against returning visitors; traffic sources;
funnel conversion from start to end of purchase & Mentions; followers; buzz, a
combined popularity measure; share of voice, conversations about you against the
competition; sentiment\\
\rlab{Major players} & Google Analytics, Adobe Analytics, Webtrends &
Facebook Insights, X Analytics, Hootsuite\\
\end{mmlongtable}

\section{The Role of Web Analytics in Digital Marketing}

\begin{keybox}[High Yield]
``Explain the role of web analytics in digital marketing. How can businesses
benefit from analysing website data?''
\pyq{CIE-II, 7 May 2025 --- 10 M}
The question has two halves; answer both. The roles are what analytics
\emph{does}; the benefits are what the business \emph{gets}.
\end{keybox}

\begin{enumerate}
\item \sfb{Measurement of effectiveness} --- establishes which campaigns,
      channels and creatives actually produce results.
\item \sfb{Understanding customer behaviour} --- reveals the real path users
      take, not the path the business assumed.
\item \sfb{Conversion optimisation} --- locates the exact step at which
      prospects abandon.
\item \sfb{Budget allocation} --- directs spend to the lowest cost per
      acquisition.
\item \sfb{Content strategy} --- identifies which topics and formats attract and
      hold the audience.
\item \sfb{Technical improvement} --- surfaces slow pages, broken journeys and
      device-specific failures.
\item \sfb{Segmentation and personalisation} --- enables different treatment of
      different audiences.
\item \sfb{Forecasting and target-setting} --- historical data makes goals
      realistic and progress measurable.
\end{enumerate}

\subsection{Benefits of analysing website data}

\begin{itemize}
\item Higher conversion rates, from evidence-based fixes rather than opinion.
\item Lower customer acquisition cost, through better channel allocation.
\item Improved user experience and reduced bounce.
\item Faster detection of problems --- a tracking break, a broken checkout, a
      traffic collapse.
\item Better content decisions and less wasted production effort.
\item Demonstrable ROI, which protects and justifies the marketing budget.
\item Competitive advantage, from acting on data faster than rivals do.
\end{itemize}

\section{Aligning Analytics KPIs With Organisational Goals}

\begin{exambox}[Model Answer --- 10 Marks]
\pyq{CIE-I, 16 Apr 2026 --- 10 M}
\begin{enumerate}
\item \sfb{Identify business goals} --- increase sales, brand awareness\ix{brand awareness}, lead
      generation.
\item \sfb{Map goals to relevant KPIs.} Sales growth \ra{} conversion rate,
      revenue, cart-abandonment rate. Brand awareness \ra{} website traffic,
      impressions, reach. Lead generation\ix{lead generation} \ra{} form submissions, sign-ups.
      Customer engagement \ra{} bounce rate, session duration, pages per visit.
\item \sfb{Use SMART\ix{SMART} KPIs} --- Specific, Measurable, Achievable, Relevant,
      Time-bound. For example: increase conversion rate by 10\% in three months.
\item \sfb{Use analytics tools} --- track the KPIs through Google Analytics\ix{Google Analytics} and
      dashboards.
\item \sfb{Monitor and optimise continuously} --- analyse data regularly and
      modify strategy on the basis of measured performance.
\end{enumerate}
\end{exambox}

\section{A/B Testing and Conversion Rate Optimisation}

\begin{definitionbox}[A/B Testing and CRO]
\keytermas{A/B testing}{A/B testing}, also called split testing, is showing two
versions of a page or element to two randomly assigned halves of the traffic
simultaneously, and measuring which produces a statistically significantly
better outcome on a defined goal.

\keytermas{Conversion rate optimisation}{conversion rate optimisation} (CRO) is
the systematic process of increasing the percentage of visitors who complete a
desired action, using analytics, user research and controlled experiments.
\end{definitionbox}

\begin{figure}[htbp]
\centering
\begin{tikzpicture}[font=\sffamily\footnotesize,
    pg/.style={draw=slate!55,line width=0.8pt,minimum width=2.5cm,
               minimum height=1.9cm,align=center,font=\sffamily\scriptsize,
               text=slate}]
  \node[draw=slate,line width=1pt,fill=white,rounded corners=3pt,
        minimum width=2.6cm,minimum height=0.85cm,align=center,
        font=\sffamily\bfseries\scriptsize,text=slate] (t) at (0,2.1)
        {INCOMING TRAFFIC};

  \node[pg,fill=ocrepale] (a) at (-2.0,0)  {\textbf{VERSION A}\\[3pt] the control\\[3pt] 50\% of traffic};
  \node[pg,fill=deepbluepale] (b) at (2.0,0) {\textbf{VERSION B}\\[3pt] the variant\\[3pt] 50\% of traffic};

  \draw[-{Latex[length=2.3mm]},ocre,line width=1pt]     (t.south) -- (-2.0,1.0);
  \draw[-{Latex[length=2.3mm]},deepblue,line width=1pt] (t.south) -- ( 2.0,1.0);

  \node[draw=greenok,line width=1pt,fill=greenpale,rounded corners=3pt,
        minimum width=6.2cm,minimum height=0.85cm,align=center,
        font=\sffamily\bfseries\scriptsize,text=greenok] (m) at (0,-1.7)
        {MEASURE ONE GOAL, RUN TO STATISTICAL SIGNIFICANCE};
  \draw[-{Latex[length=2.3mm]},slate,line width=0.9pt] (-2.0,-1.0) -- (-1.4,-1.28);
  \draw[-{Latex[length=2.3mm]},slate,line width=0.9pt] ( 2.0,-1.0) -- ( 1.4,-1.28);

  \node[anchor=west,align=left,text width=4.6cm,font=\sffamily\scriptsize,
        text=slate] at (3.7,0.4)
    {\textbf{One variable at a time.}\\[3pt]
     Changing two elements at once means a win cannot be attributed to either
     of them.};
  \node[anchor=west,align=left,text width=4.6cm,font=\sffamily\scriptsize,
        text=ocredark] at (3.7,-1.7)
    {A 10\% lift applies to \emph{all} future traffic at no extra acquisition
     cost --- which is why CRO compounds.};
\end{tikzpicture}
\caption[The A/B testing mechanism]%
        {A/B testing splits traffic simultaneously rather than sequentially, so
         that seasonality, campaign timing and traffic-mix changes affect both
         versions equally.}
\label{fig:abtest}
\end{figure}

\begin{itemize}
\item \sfb{Significance.} It replaces opinion and seniority with evidence;
      isolates the effect of a single variable; produces compounding gains, since
      a 10\% lift applies to all future traffic at zero additional acquisition
      cost; reduces the risk of large redesigns by testing incrementally; and
      reveals genuine user preference, which frequently contradicts internal
      assumption.
\item \sfb{What to test.} Headlines; calls to action --- text, colour, position;
      hero images; form length and field order; page layout; pricing
      presentation; trust signals; and the checkout flow.
\item \sfb{Process.} Analyse data to find the weakest step \ra{} form a
      hypothesis (``shortening the form from nine fields to four will raise
      submissions'') \ra{} build the variant \ra{} split traffic \ra{} run to
      statistical significance \ra{} implement the winner \ra{} repeat.
\item \sfb{Tools.} VWO, Optimizely and AB Tasty for the experiments, plus Hotjar\ix{Hotjar}
      and Microsoft Clarity for the qualitative diagnosis that generates the
      hypotheses in the first place.
\end{itemize}

\begin{summarybox}
Web analytics collects, measures, analyses and reports website data to understand
and optimise web usage; it is what makes digital marketing measurable at all.
Against traditional analysis it wins on speed, accuracy, cost, sample size,
attribution and the continuity of its feedback loop. Metrics fall into three
families --- volume, quality and action --- and no family can be read alone.
Bounce rate is a single-page no-interaction session and must be read in context;
exit rate is merely the last page of any session. Dimensions are attributes and
form the rows; metrics are numbers and form the columns. Web analytics plays
eight roles from measuring effectiveness to forecasting, and yields seven
business benefits from higher conversion to demonstrable ROI. KPIs are aligned to
goals in five steps. A/B testing and CRO replace opinion with evidence and
compound, because a lift applies to all future traffic at no extra acquisition
cost.
\end{summarybox}

%%%%%%%%%%%%%%%%%%%%%%%%%%%%%%%%%%%%%%%%%%%%%%%%%%%%%%%%%%%%%%%%%%%%%%%%%%%%%%
\chapter{Google Analytics Tools}
%%%%%%%%%%%%%%%%%%%%%%%%%%%%%%%%%%%%%%%%%%%%%%%%%%%%%%%%%%%%%%%%%%%%%%%%%%%%%%

\syllabusstrip{Google Analytics Tools.}

\begin{learningoutcomes}
\begin{itemize}
\item The purpose of Google Analytics.
\item Its key features --- the four to name, and the full list.
\item Acquisition channels, and the optimisation action for each.
\item The seven report families.
\item The eight-step process for setting up and tracking goals.
\item Google Tag Manager\ix{Google Tag Manager} --- its contribution and its governance risks.
\end{itemize}
\end{learningoutcomes}

\section{Purpose}

\begin{exambox}[Model Answer --- 2 Marks]
\pyq{CIE-II, 7 May 2025 --- 2 M}
The main purpose of \keytermas{Google Analytics}{Google Analytics} is to collect,
measure and report data about website and app users --- who they are, how they
arrived, what they did and whether they converted --- so that businesses can
understand user behaviour, evaluate marketing performance and make
evidence-based decisions to improve traffic quality and conversions.
\end{exambox}

\section{Key Features}

\begin{exambox}[Four Features --- 2 Marks]
\pyq{CIE-II, 26 May 2026 --- 2 M}
(1) Traffic source analysis, (2) bounce rate measurement, (3) user demographics
tracking, (4) conversion and goal tracking.
\end{exambox}

\noindent A fuller list, for eight- and ten-mark answers:

\begin{itemize}
\item \sfb{Real-time reporting} --- who is on the site right now, and what they
      are doing.
\item \sfb{Audience reports} --- demographics, interests, geography, language,
      technology and device.
\item \sfb{Acquisition reports} --- which channels, sources, mediums and
      campaigns delivered the traffic.
\item \sfb{Behaviour and engagement reports} --- landing pages, exit pages, site
      content performance, site search, page load times.
\item \sfb{Conversion and goal tracking} --- destination, duration,
      pages-per-session and event goals; assigning monetary value to conversions.
\item \sfb{E-commerce tracking} --- product performance, shopping and checkout
      behaviour, average order value, revenue.
\item \sfb{Funnel visualisation} --- where users drop out of a defined
      multi-step path.
\item \sfb{Segments} --- isolating a subset of users, for example mobile users
      from paid search who did not convert.
\item \sfb{Attribution and model comparison} --- how credit for a conversion is
      distributed across touchpoints.
\item \sfb{Custom dashboards, reports and alerts.}
\item \sfb{Audience building and remarketing lists} exported to the advertising
      platform.
\item \sfb{Integrations} --- Google Ads, Search Console, BigQuery, Google Tag
      Manager, Looker Studio.
\end{itemize}

\section{Acquisition Channels}

\begin{keybox}[High Yield]
``Discuss the concept of `Acquisition channels' in Google Analytics. How can
business use this information to optimize their marketing efforts?''
\pyq{SEE, Oct/Nov 2023, Q5a --- 8 M}
\end{keybox}

\begin{mmlongtable}{Acquisition channels and their optimisation actions}{tab:channels3}%
{@{}P{0.17\textwidth}P{0.37\textwidth}P{0.39\textwidth}@{}}
\rowcolor{ocre}\thd{Channel} & \thd{What it means} & \thd{Typical optimisation action}\\
\midrule
\endfirsthead
\rowcolor{ocre}\thd{Channel} & \thd{What it means} & \thd{Typical optimisation action}\\
\midrule
\endhead
\rlab{Organic search} & Unpaid clicks from search engines & Invest in SEO for the
pages and queries already converting\\
\rowcolor{tablealt}
\rlab{Paid search} & Clicks from search advertisements & Shift budget to keywords
with the lowest cost per conversion\\
\rlab{Direct} & Users typing the URL or using bookmarks --- also the catch-all for
untagged traffic & Build brand; audit for missing UTM tags inflating this
bucket\\
\rowcolor{tablealt}
\rlab{Referral} & Links from other websites & Nurture the partnerships and
publications actually sending buyers\\
\rlab{Organic social} & Unpaid social posts & Double down on the platform
delivering engaged traffic\\
\rowcolor{tablealt}
\rlab{Paid social} & Social advertising & Reallocate spend by cost per result\\
\rlab{E-mail} & Clicks from e-mail campaigns & Increase send frequency to the
highest-converting segments\\
\rowcolor{tablealt}
\rlab{Display} & Banner and remarketing advertisements & Refine placements and
exclude low-quality inventory\\
\rlab{Affiliates} & Partner-driven traffic & Adjust commission by partner
quality\\
\end{mmlongtable}

\begin{exambox}[How the Business Uses This]
It identifies which channels deliver not just volume but \emph{converting}
volume; it exposes the cost per acquisition of each channel so that budget can be
reallocated; it reveals \emph{assisted conversions}, where a channel contributes
without taking the last click; it enables channel-specific landing page
optimisation; and it prevents the classic error of judging a channel by traffic
when it should be judged by revenue.
\end{exambox}

\begin{warnbox}[An Inflated ``Direct'' Bucket Is a Measurement Fault]
Direct traffic should be people who typed the URL or used a bookmark. In practice
it is also where every untagged link ends up --- an e-mail without UTM
parameters, a link in a PDF, a QR code. A large and growing Direct channel is
therefore more often a tagging failure than a brand success, and it silently
steals credit from the channels that earned it.
\end{warnbox}

\section{Report Families}

\begin{keybox}[High Yield]
``List and explain the various reports that can be generated using Google
Analytics.''
\pyq{SEE, Oct/Nov 2023, Q6b --- 8 M}
\end{keybox}

\begin{mmlongtable}{The report families in Google Analytics}{tab:reports}%
{@{}P{0.24\textwidth}P{0.69\textwidth}@{}}
\rowcolor{ocre}\thd{Report family} & \thd{What it tells you}\\
\midrule
\endfirsthead
\rowcolor{ocre}\thd{Report family} & \thd{What it tells you}\\
\midrule
\endhead
\rlab{Real-time} & Active users right now, their pages, sources and locations ---
used to verify a campaign launch or a tag deployment\\
\rowcolor{tablealt}
\rlab{Audience / User} & Who the visitors are: demographics, interests,
geography, language, technology, devices, new against returning, user lifetime
value, cohort retention\\
\rlab{Acquisition / Traffic acquisition} & Where users came from --- channel,
source and medium, campaign, plus advertising and Search Console integration\\
\rowcolor{tablealt}
\rlab{Behaviour / Engagement} & What users did on the site --- pages and screens,
landing pages, exit pages, site search, page timings, events\\
\rlab{Conversions / Monetisation} & Goal completions, goal value, funnel
visualisation, e-commerce revenue, product performance, checkout behaviour,
purchase journey\\
\rowcolor{tablealt}
\rlab{Attribution / Advertising} & How conversion credit is distributed across
touchpoints; model comparison\\
\rlab{Custom reports and Explorations} & User-defined combinations of dimensions
and metrics; funnel explorations, path explorations, segment overlap\\
\end{mmlongtable}

\section{Setting Up and Tracking Goals}

\begin{keybox}[High Yield]
``Describe the process of setting up and tracking goals in Google Analytics.''
\pyq{SEE, Oct/Nov 2023, Q10b --- 6 M}
\end{keybox}

\begin{enumerate}
\item \sfb{Define what a conversion means} for the business --- a purchase, a
      form submission, a newsletter sign-up, a brochure download, a call click.
\item \sfb{Choose the goal type} --- destination (a thank-you URL is reached),
      duration (a session exceeds a set time), pages or screens per session, or
      event (a specific interaction such as a video play or a button click). In
      GA4 these are configured as events marked as conversions.
\item \sfb{Implement the tracking} --- deploy the tag via Google Tag Manager,
      define the trigger, and fire the event on the qualifying action.
\item \sfb{Assign a monetary value} to the goal where possible, so that
      non-transactional conversions can still be compared and optimised.
\item \sfb{Configure the funnel} --- define the required steps preceding the
      goal, so that drop-off is visible step by step.
\item \sfb{Verify} using real-time reports and debug mode \emph{before} relying
      on the data.
\item \sfb{Track and report} --- monitor conversion rate by channel, campaign,
      device and landing page.
\item \sfb{Act on the findings} --- fix the step with the largest drop-off,
      reallocate budget to the highest-converting channels, and A/B test the
      weakest page.
\end{enumerate}

\begin{keybox}[Why Goal Tracking Matters]
Without goals, analytics measures \emph{activity} rather than \emph{outcome}.
Goals are what convert traffic reports into ROI reports.
\end{keybox}

\section{Google Tag Manager}

\begin{keybox}[High Yield]
``Critically evaluate the contribution of Google Tag Manager in optimizing
marketing channels for global companies.''
\pyq{SEE, Jun/Jul 2025, Q6b --- 8 M}
The word \emph{critically} requires both halves: the contribution \emph{and} the
limitations. An answer giving only the benefits is incomplete.
\end{keybox}

\begin{definitionbox}[Google Tag Manager]
\keytermas{Google Tag Manager}{Google Tag Manager} is a container that lets
marketers deploy, edit and version tracking tags without editing site code.
\end{definitionbox}

\subsection{The contribution}

\begin{itemize}
\item It lets marketers deploy, edit and version tracking tags without a
      developer, which collapses the change cycle from weeks to minutes.
\item It standardises the data layer across dozens of regional sites, so that
      global reporting is comparable.
\item It centralises consent management, which is essential under GDPR and
      similar regimes.
\item It reduces page weight by loading tags asynchronously.
\item It enables rapid A/B and campaign experimentation.
\item It provides preview and version rollback, reducing the risk of breaking
      measurement.
\end{itemize}

\subsection{The critical view}

\begin{warnbox}[Where Tag Manager Becomes the Problem]
It can become a governance problem: untracked tag sprawl, duplicate conversion
firing, degraded page performance, and security exposure if publish rights are
granted too widely.

\smallskip
It also does not fix a poor measurement plan. It only makes deploying one
faster.
\end{warnbox}

\begin{summarybox}
Google Analytics collects, measures and reports data about users --- who they
are, how they arrived, what they did and whether they converted. The four
features to name are traffic source analysis, bounce rate measurement,
demographics tracking, and conversion and goal tracking; the fuller list adds
real-time, audience, acquisition, behaviour, e-commerce, funnels, segments,
attribution, dashboards, remarketing audiences and integrations. Acquisition
channels group traffic by origin, each with its own optimisation action, and the
purpose of the report is to judge channels by revenue rather than volume. Reports
fall into seven families. Goals are configured in eight steps, verified before
being trusted, and are what turn traffic reports into ROI reports. Tag Manager
collapses the deployment cycle and standardises global measurement, but
introduces governance risk and cannot substitute for a measurement plan.
\end{summarybox}

%%%%%%%%%%%%%%%%%%%%%%%%%%%%%%%%%%%%%%%%%%%%%%%%%%%%%%%%%%%%%%%%%%%%%%%%%%%%%%
\chapter{Web Analytics Tools}
%%%%%%%%%%%%%%%%%%%%%%%%%%%%%%%%%%%%%%%%%%%%%%%%%%%%%%%%%%%%%%%%%%%%%%%%%%%%%%

\syllabusstrip{Web Analytics Tools.}

\begin{learningoutcomes}
\begin{itemize}
\item The main web analytics tools other than Google Analytics, and what each is
      good at.
\item How quantitative and qualitative tools combine --- the ``what'' and the
      ``why''.
\end{itemize}
\end{learningoutcomes}

\section{Tools Other Than Google Analytics}

\begin{exambox}[Two Tools --- 2 Marks]
``Name any two popular web analytics tools other than Google Analytics.''
\pyq{CIE-II, 7 May 2025 --- 2 M}
\textbf{Adobe Analytics} and \textbf{Matomo}. Also acceptable: Webtrends,
Mixpanel, Hotjar, Microsoft Clarity, Kissmetrics, Piwik PRO, Heap, Amplitude,
Crazy Egg, SimilarWeb.
\end{exambox}

\begin{mmlongtable}{Web analytics tools and their strengths}{tab:watools}%
{@{}P{0.26\textwidth}P{0.67\textwidth}@{}}
\rowcolor{ocre}\thd{Tool} & \thd{Strength}\\
\midrule
\endfirsthead
\rowcolor{ocre}\thd{Tool} & \thd{Strength}\\
\midrule
\endhead
\rlab{Adobe Analytics} & Enterprise-grade segmentation, attribution and
integration with the Adobe Experience Cloud\\
\rowcolor{tablealt}
\rlab{Matomo} \newline {\scriptsize formerly Piwik} & Open-source and
self-hostable; full data ownership and privacy compliance\\
\rlab{Webtrends} & A long-established enterprise web analytics platform\\
\rowcolor{tablealt}
\rlab{Mixpanel and Amplitude} & Product analytics --- event-based funnels,
cohorts and retention for apps and software services\\
\rlab{Hotjar, Microsoft Clarity, Crazy Egg} & Qualitative behaviour ---
heatmaps, scroll maps, click maps and session recordings\\
\rowcolor{tablealt}
\rlab{Kissmetrics} & Person-level tracking across sessions and devices\\
\rlab{SimilarWeb} & Competitive traffic estimation and benchmarking\\
\rowcolor{tablealt}
\rlab{Google Search Console} & Search-side data --- impressions, clicks, average
position, indexing and Core Web Vitals\\
\rlab{Google Tag Manager} & Central deployment and version control of all
tracking tags without editing site code\\
\rowcolor{tablealt}
\rlab{Looker Studio} & Dashboarding and reporting across multiple data sources\\
\end{mmlongtable}

\section{The Quantitative and the Qualitative}

\begin{keybox}[High Yield]
``Demonstrate how tools like Google Analytics and Hotjar can be used to develop
insights into user behavior and improve marketing decisions.''
\pyq{SEE, Jun/Jul 2025, Q7a --- 8 M}
The organising idea for the answer: analytics gives the \textbf{what},
behavioural tools give the \textbf{why}, and neither is sufficient alone.
\end{keybox}

\begin{figure}[htbp]
\centering
\begin{tikzpicture}[font=\sffamily\footnotesize,
    side/.style={rounded corners=4pt,line width=1pt,align=left,
                 text width=4.85cm,inner sep=8pt,font=\sffamily\scriptsize,
                 text=slate}]
  \node[side,draw=deepblue,fill=deepbluepale] (q) at (0,0)
    {\textbf{\textcolor{deepblue}{GOOGLE ANALYTICS --- the \emph{what}}}\\[4pt]
     which channels deliver traffic\\
     which landing pages bounce\\
     where users drop out of checkout\\
     which segments convert\\
     what each channel costs per acquisition\\[4pt]
     \textit{quantitative, complete, anonymous}};
  \node[side,draw=ocre,fill=ocrepale] (h) at (5.55,0)
    {\textbf{\textcolor{ocredark}{HOTJAR --- the \emph{why}}}\\[4pt]
     heatmaps: what attracts attention\\
     scroll maps: is the CTA ever reached\\
     click maps: rage clicks on dead elements\\
     recordings: hesitation and confusion\\
     on-page surveys: the reason, in their words\\[4pt]
     \textit{qualitative, sampled, observed}};

  % text width is required as well as align=center: without it the node never
  % wraps and grows to the width of its longest unbroken line.
  \node[draw=greenok,line width=1pt,fill=greenpale,rounded corners=4pt,
        minimum width=10.4cm,text width=10.1cm,align=center,
        font=\sffamily\scriptsize,text=slate,inner sep=6pt] at (2.78,-2.55)
    {\textbf{\textcolor{greenok}{THE COMBINED WORKFLOW}}\\[3pt]
     analytics shows the pricing page loses 60\% of visitors \ra{} recordings
     reveal users clicking a non-clickable comparison table \ra{} the table is
     made interactive \ra{} an A/B test validates the change \ra{} analytics
     confirms the improved funnel conversion};
\end{tikzpicture}
\caption[Quantitative and qualitative analytics combined]%
        {Analytics locates the problem; behavioural tools explain it. A funnel
         report can say \emph{where} users leave but never \emph{why}, and a
         session recording can show why but never how often.}
\label{fig:whatwhy}
\end{figure}

\begin{exambox}[Marketing Decisions Improved]
Budget reallocation, landing page redesign priority, form simplification, content
placement and mobile-specific fixes --- each grounded in evidence rather than
opinion.
\end{exambox}

\begin{summarybox}
Beyond Google Analytics, the toolset divides by job: Adobe Analytics for
enterprise attribution, Matomo for self-hosted data ownership, Mixpanel and
Amplitude for product event analytics, Hotjar and Clarity for qualitative
behaviour, Kissmetrics for person-level tracking, SimilarWeb for competitive
benchmarking, Search Console for search-side data, Tag Manager for deployment and
Looker Studio for reporting. Quantitative tools answer what happened and how
often; qualitative tools answer why. The productive workflow alternates between
them: analytics locates a drop-off, recordings explain it, a test validates the
fix, and analytics confirms the result.
\end{summarybox}

%%%%%%%%%%%%%%%%%%%%%%%%%%%%%%%%%%%%%%%%%%%%%%%%%%%%%%%%%%%%%%%%%%%%%%%%%%%%%%
%  Unit III, second half: E-mail Marketing (Chapters 4-7)
%%%%%%%%%%%%%%%%%%%%%%%%%%%%%%%%%%%%%%%%%%%%%%%%%%%%%%%%%%%%%%%%%%%%%%%%%%%%%%

\chapter{Effective E-mail Campaigns}

\syllabusstrip{E-mail Marketing: Effective E-mail Campaigns.}

\begin{learningoutcomes}
\begin{itemize}
\item What e-mail marketing is, and why it is an owned channel.
\item E-mail 1.0 against E-mail 2.0.
\item The advantages, and the liabilities that come with them.
\item Opt-in against spam, and the legal position.
\item The anatomy of an effective e-mail --- three rules and a checklist.
\item The content arsenal: nine e-mail types in three strategic groups.
\item Transactional against promotional e-mail.
\end{itemize}
\end{learningoutcomes}

\section{What E-mail Marketing Is}

\begin{definitionbox}[E-mail Marketing]
\keytermas{E-mail marketing}{e-mail marketing} is the use of e-mail to deliver
commercial and relationship-building messages to a permission-based list of
subscribers, in order to acquire, engage, convert and retain customers.
\end{definitionbox}

It is a vital tool used throughout the entire sales funnel\ix{sales funnel} --- engaging
prospects, maintaining interest and driving key revenue for online retailers at a
fraction of the cost of direct mail.

\begin{keybox}[Why E-mail Is Structurally Different]
Every other channel rents its audience. Social reach is granted by an algorithm
that can withdraw it; search ranking is granted by an algorithm that can change.
An e-mail list is an \textbf{owned asset}: no intermediary stands between the
message and the recipient, which is why e-mail has the highest conversion and
retention rate of any digital channel.
\end{keybox}

\section{E-mail 1.0 Against E-mail 2.0}

\begin{mmlongtable}{The generational shift in e-mail marketing}{tab:email12}%
{@{}P{0.44\textwidth}P{0.49\textwidth}@{}}
\rowcolor{ocre}\thd{E-mail 1.0} & \thd{E-mail 2.0}\\
\midrule
\endfirsthead
\rowcolor{ocre}\thd{E-mail 1.0} & \thd{E-mail 2.0}\\
\midrule
\endhead
Simple one-way messages & Behaviour-based triggers\\
\rowcolor{tablealt}
Static auto-responders & Segmented targeting based on purchase history and
interests\\
Mass ``batch and blast'' distribution & Dynamic lifecycle automation\\
\rowcolor{tablealt}
Untracked engagement & Deep tracking metrics\\
\end{mmlongtable}

\section{Advantages}

\begin{exambox}[Four Advantages Over Traditional Methods --- 2 Marks]
\pyq{CIE-II, 26 May 2026 --- 2 M}
Low cost; global reach; easy performance tracking; personalised communication.
\end{exambox}

\begin{itemize}
\item \sfb{Cost and speed} --- cheaper and faster than traditional direct mail.
\item \sfb{Tracking} --- real-time visibility into opens, clicks and conversions.
\item \sfb{Personalisation} --- dynamic tailoring based on collected data.
\item \sfb{Global reach} at no incremental cost per recipient.
\item \sfb{Ownership of the channel} --- unlike social media, the list is an
      owned asset that no algorithm can throttle.
\item \sfb{Highest conversion and retention rate} of any digital channel.
\item \sfb{Direct and permission-based} --- the recipient asked to hear from you.
\end{itemize}

\section{Disadvantages and Liabilities}

\begin{itemize}
\item \sfb{Spam and legal risk} --- filters block messages, and anti-spam
      legislation carries severe penalties.
\item \sfb{Engagement decay} --- constant tweaking is required to prevent list
      fatigue.
\item \sfb{Code rendering} --- appearance breaks across varying e-mail clients,
      servers and settings, for example when images are blocked by default.
\item \sfb{Hidden costs} --- scaling campaigns requires design, management and
      software investment, potentially upward of \$1,000 per month.
\item \sfb{Deliverability management} --- sender reputation, authentication (SPF,
      DKIM, DMARC) and list hygiene require ongoing attention.
\item \sfb{Inbox competition} --- the average inbox is crowded, so open rates are
      structurally limited.
\end{itemize}

\section{Opt-In Against Spam}

\begin{warnbox}[The Legal Point]
The CAN-SPAM Act levies fines of up to \textbf{\$16,000 per violation} for
non-compliance. Permission is not a courtesy; it is a legal requirement. Indian
and EU equivalents --- the IT Act rules and GDPR --- impose comparable
obligations.
\end{warnbox}

\begin{mmlongtable}{Spam against opt-in}{tab:optin}%
{@{}P{0.17\textwidth}P{0.38\textwidth}P{0.38\textwidth}@{}}
\rowcolor{ocre}\thd{} & \thd{Spam} & \thd{Opt-In}\\
\midrule
\endfirsthead
\rowcolor{ocre}\thd{} & \thd{Spam} & \thd{Opt-In}\\
\midrule
\endhead
\rlab{Acquisition method} & Unsolicited; default web checkboxes; purchased lists &
Willingly provided via lead forms in exchange for incentives\\
\rowcolor{tablealt}
\rlab{Expectation setting} & Blind blasting; unpredictable frequency & Clear
promises about e-mail type --- sales against tips --- and frequency\\
\rlab{Brand impact} & Degrades reputation; perceived as a nuisance & Builds
long-term trust; establishes authority; boosts upsells\\
\rowcolor{tablealt}
\rlab{List quality} & High volume, low intent & Micro-focused, highly targeted
engagement\\
\end{mmlongtable}

\begin{definitionbox}[Single and Double Opt-In]
\keytermas{Single opt-in}{single opt-in} adds the address immediately on
submission. \keytermas{Double opt-in}{double opt-in} requires the subscriber to
confirm through a verification e-mail --- producing a smaller but far cleaner and
more engaged list, improving deliverability and providing documentary proof of
consent.
\end{definitionbox}

\begin{exambox}[Why Segmentation Matters in E-mail --- 2 Marks]
\pyq{SEE, Jun/Jul 2025 --- 2 M}
Segmentation divides the list by demographics, purchase history, behaviour,
engagement level or lifecycle stage so that each group receives relevant content.
It raises open and click rates, reduces unsubscribes and spam complaints,
improves deliverability by maintaining engagement, and enables personalisation
--- which is the single largest driver of e-mail conversion.
\end{exambox}

\section{Anatomy of an Effective E-mail}

\subsection{The three rules}

\begin{frameworkbox}[The Three Rules]
\textbf{1. One topic.} Restrict the e-mail to a single marketing message.
Multiple messages create a ``busy'' user experience; focus drives higher revenue
per message.

\textbf{2. Attractive, simple design.} The design must complement the message,
not distract from it. Design on the assumption that mobile clients may disable
images.

\textbf{3. Clear call to action\ix{call to action}.} The reader must never wonder what to do next;
the path to purchase must be frictionless.
\end{frameworkbox}

\begin{keybox}[The Result of Experience]
Continuously A/B test designs and offers in order to learn customer behaviour. No
template is permanently correct.
\end{keybox}

\subsection{Checklist for an effective marketing e-mail}

\begin{itemize}
\item A \sfb{subject line} that is specific, benefit-led and under about 50
      characters --- this alone determines the open rate.
\item A \sfb{recognisable, consistent sender name} --- people open based on who
      sent it as much as on what it says.
\item \sfb{Preheader text} that extends the subject line rather than repeating
      it.
\item \sfb{Personalisation beyond the first name} --- reference actual behaviour
      or past purchases.
\item \sfb{Above-the-fold value} --- the offer must be visible without scrolling.
\item A \sfb{single, prominent, repeated CTA button} with active verb text.
\item \sfb{Responsive, mobile-first design} --- the majority of opens are on
      mobile.
\item \sfb{Alt text on every image}, so the message survives image blocking.
\item A \sfb{plain-text fallback} version.
\item A \sfb{visible unsubscribe link} and the sender's physical address --- both
      legally required.
\item \sfb{Tested rendering} across major clients before sending.
\end{itemize}

\section{The Content Arsenal --- Nine Types of Marketing E-mail}

\begin{keybox}[Name the Group, Not Just the Type]
The nine templates fall into three strategic groups. Naming the group and the
rules for each type is what separates a full-mark answer from a list.
\end{keybox}

\begin{figure}[htbp]
\centering
\begin{tikzpicture}[font=\sffamily\footnotesize,
    grp/.style={rounded corners=4pt,line width=1pt,align=center,
                text width=3.35cm,inner sep=7pt,font=\sffamily\scriptsize,
                text=slate}]
  \node[grp,draw=warnred,fill=warnpale] at (0,0)
    {\textbf{\textcolor{warnred}{GROUP 1}}\\ \textbf{DRIVE REVENUE}\\[5pt]
     Promotional\\ New Inventory\\ Reorder};
  \node[grp,draw=ocre,fill=ocrepale] at (4.05,0)
    {\textbf{\textcolor{ocredark}{GROUP 2}}\\ \textbf{BUILD RELATIONSHIPS}\\[5pt]
     Newsletter\\ Welcome\\ Educational\\ Product Advice};
  \node[grp,draw=deepblue,fill=deepbluepale] at (8.1,0)
    {\textbf{\textcolor{deepblue}{GROUP 3}}\\ \textbf{GATHER INTELLIGENCE}\\[5pt]
     Testimonials\\ Surveys};
\end{tikzpicture}
\caption[The nine e-mail types in three strategic groups]%
        {The content arsenal. The group determines the tone and the metric: Group~1
         is judged on revenue, Group~2 on engagement, Group~3 on response rate.}
\label{fig:ninetypes}
\end{figure}

\subsection{Group 1 --- drive revenue}

\begin{mmlongtable}{Group 1: revenue e-mails}{tab:group1}%
{@{}P{0.15\textwidth}P{0.28\textwidth}P{0.50\textwidth}@{}}
\rowcolor{ocre}\thd{Type} & \thd{Goal} & \thd{Rules}\\
\midrule
\endfirsthead
\rowcolor{ocre}\thd{Type} & \thd{Goal} & \thd{Rules}\\
\midrule
\endhead
\rlab{Promotional} & Entice immediate purchase & Keep it short; make the offer
front and centre; create a sense of urgency with active language\\
\rowcolor{tablealt}
\rlab{New Inventory} & Alert customers to new arrivals & Send immediately upon
arrival; prioritise a strong product photograph; minimise text; tease the reveal
in the subject line\\
\rlab{Reorder} & Replenish consumables & Make the buying process simple with a
``Reorder Now'' button; remind them of the value and the hassle avoided;
explicitly mention past purchases\\
\end{mmlongtable}

\subsection{Group 2 --- build relationships}

\begin{mmlongtable}{Group 2: relationship e-mails}{tab:group2}%
{@{}P{0.15\textwidth}P{0.28\textwidth}P{0.50\textwidth}@{}}
\rowcolor{ocre}\thd{Type} & \thd{Goal} & \thd{Rules}\\
\midrule
\endfirsthead
\rowcolor{ocre}\thd{Type} & \thd{Goal} & \thd{Rules}\\
\midrule
\endhead
\rlab{Newsletter} & Improve brand awareness & Think ``mini-newspaper''; break copy
into digestible sections; use simple layouts; always include contact
information\\
\rowcolor{tablealt}
\rlab{Welcome} & Establish a relationship with new opt-ins & Use a warm,
conversational tone --- a ``virtual handshake''; offer a reward such as 10\% off;
remind them of the benefits of joining\\
\rlab{Educational} & Build trust through industry knowledge & Offer relevant,
bite-sized information linking out to full blog posts; include a subtle
mini-promotion\\
\rowcolor{tablealt}
\rlab{Product Advice} & Establish authority & Solve customer problems; teach
product features; emphasise customer service; proofread meticulously to protect
credibility\\
\end{mmlongtable}

\subsection{Group 3 --- gather intelligence}

\begin{mmlongtable}{Group 3: intelligence e-mails}{tab:group3}%
{@{}P{0.15\textwidth}P{0.28\textwidth}P{0.50\textwidth}@{}}
\rowcolor{ocre}\thd{Type} & \thd{Goal} & \thd{Rules}\\
\midrule
\endfirsthead
\rowcolor{ocre}\thd{Type} & \thd{Goal} & \thd{Rules}\\
\midrule
\endhead
\rlab{Testimonials} & Reinforce business value through social proof & Use a sleek
design to frame quotes; pair text with powerful images --- customer photographs
or product shots; offer a link for readers to leave their own feedback\\
\rowcolor{tablealt}
\rlab{Surveys} & Collect data to improve customer experience & Clearly explain
the incentive, such as a prize draw; explicitly address why the data is being
collected; provide a highly visible, clickable access button\\
\end{mmlongtable}

\section{Transactional Against Promotional E-mail}

\begin{keybox}[High Yield]
``Describe the difference between transactional emails and promotional emails.
Provide examples of each and explain when to use them.''
\pyq{SEE, Oct/Nov 2023, Q6a --- 8 M}
\end{keybox}

\begin{mmlongtable}{Transactional against promotional e-mail}{tab:transvspromo}%
{@{}P{0.14\textwidth}P{0.40\textwidth}P{0.39\textwidth}@{}}
\rowcolor{ocre}\thd{Basis} & \thd{Transactional E-mail} & \thd{Promotional E-mail}\\
\midrule
\endfirsthead
\rowcolor{ocre}\thd{Basis} & \thd{Transactional E-mail} & \thd{Promotional E-mail}\\
\midrule
\endhead
\rlab{Trigger} & A specific action or event performed by the individual user &
A marketing calendar or campaign decision\\
\rowcolor{tablealt}
\rlab{Purpose} & To deliver information the user expects and needs & To generate
awareness, demand or sales\\
\rlab{Consent required} & Implied by the transaction itself & Explicit opt-in
required\\
\rowcolor{tablealt}
\rlab{Volume} & One-to-one, automated & One-to-many, batched or segmented\\
\rlab{Open rate} & Very high, often 40--60\% or more, because it is expected &
Much lower, typically 15--25\%\\
\rowcolor{tablealt}
\rlab{Examples} & Order confirmation, shipping and delivery notification,
password reset, one-time password, invoice or receipt, account creation,
appointment reminder & Sale announcement, newsletter, new-product launch,
seasonal offer, webinar invitation, re-engagement campaign\\
\rlab{When to use} & Immediately after the triggering event --- speed and clarity
matter more than design & At planned intervals against a content calendar,
targeted to segments\\
\end{mmlongtable}

\begin{keybox}[A Practical Point Worth Marks]
Because transactional e-mails are opened far more often, a small, tasteful
promotional element --- a related-product strip or a loyalty reminder ---
\emph{inside} them often outperforms an entire promotional campaign.

\smallskip
But the transactional content must remain dominant, or the message loses its
legal status as transactional and becomes subject to opt-in rules.
\end{keybox}

\begin{summarybox}
E-mail marketing delivers commercial and relationship messages to a
permission-based list in order to acquire, engage, convert and retain. It is the
only major owned channel, which is why it has the highest conversion and
retention rate. E-mail 2.0 replaces one-way blasts and static auto-responders
with behavioural triggers, segmentation and lifecycle automation. Its advantages
are cost, speed, tracking, personalisation, reach and ownership; its liabilities
are legal risk, engagement decay, rendering inconsistency, hidden cost,
deliverability management and inbox competition. Opt-in differs from spam in
acquisition, expectation, brand impact and list quality; double opt-in produces a
smaller and cleaner list. An effective e-mail keeps to one topic, a simple design
and one clear call to action. The nine templates group into revenue,
relationship and intelligence. Transactional e-mail is user-triggered, expected
and opened two to three times as often as promotional e-mail, which is
calendar-driven and requires explicit consent.
\end{summarybox}

%%%%%%%%%%%%%%%%%%%%%%%%%%%%%%%%%%%%%%%%%%%%%%%%%%%%%%%%%%%%%%%%%%%%%%%%%%%%%%
\chapter{The E-mail Plan}
%%%%%%%%%%%%%%%%%%%%%%%%%%%%%%%%%%%%%%%%%%%%%%%%%%%%%%%%%%%%%%%%%%%%%%%%%%%%%%

\syllabusstrip{E-mail Plan.}

\begin{learningoutcomes}
\begin{itemize}
\item The seven-node e-mail marketing plan, and why it is a loop.
\item A worked plan for a stated scenario.
\item The eight-step guide to segmentation and personalisation.
\end{itemize}
\end{learningoutcomes}

\section{The Seven Nodes}

\begin{keybox}[The Framework for Every ``Design an E-mail Plan'' Question]
There have been several such questions, and each is answered by walking these
seven nodes against the specific business given.
\pyq{CIE-II, 26 May 2026 --- 10 M}
\end{keybox}

\begin{figure}[htbp]
\centering
\begin{tikzpicture}[font=\sffamily\footnotesize,
    nd/.style={draw=ocre,line width=0.9pt,fill=ocrepale,circle,
               minimum size=1.72cm,align=center,
               font=\sffamily\bfseries\scriptsize,text=slate,inner sep=1pt}]
  \node[nd] (n1) at (0,0)     {1\\ choose\\ software};
  \node[nd] (n2) at (1.95,0)  {2\\ build\\ the list};
  \node[nd] (n3) at (3.90,0)  {3\\ capture\\ forms};
  \node[nd] (n4) at (5.85,0)  {4\\ define\\ goals};
  \node[nd] (n5) at (7.80,0)  {5\\ auto-\\ responders};
  \node[nd] (n6) at (9.75,0)  {6\\ add\\ triggers};
  \node[nd,draw=greenok,fill=greenpale] (n7) at (11.70,0) {7\\ monitor\\ metrics};
  \foreach \a/\b in {n1/n2,n2/n3,n3/n4,n4/n5,n5/n6,n6/n7}{
    \draw[-{Latex[length=2mm]},ocre,line width=0.9pt] (\a)--(\b);}
  \draw[-{Latex[length=2.2mm]},greenok,line width=0.9pt,dashed]
    (n7.south) -- ++(0,-1.05) -| (n1.south);
  \node[anchor=north,text=greenok,font=\sffamily\scriptsize\itshape]
    at (5.85,-1.1)
    {the metrics tell you whether the software, list, goals, sequences and
     triggers were right};
\end{tikzpicture}
\caption[The seven-node e-mail marketing plan]%
        {The seven nodes form a loop, not a line. Node~7 feeds back into Node~1.}
\label{fig:sevennodes}
\end{figure}

\begin{mmlongtable}{The seven nodes in detail}{tab:sevennodes}%
{@{}P{0.20\textwidth}P{0.73\textwidth}@{}}
\rowcolor{ocre}\thd{Node} & \thd{Action}\\
\midrule
\endfirsthead
\rowcolor{ocre}\thd{Node} & \thd{Action}\\
\midrule
\endhead
\rlab{1. Choose software} & Select a CRM or e-mail service provider based on the
capabilities required --- a simple database sender against full behavioural
campaign automation\\
\rowcolor{tablealt}
\rlab{2. Build the list} & Start capturing addresses; commit to quality over
quantity; limit sending to a maximum of about two e-mails per month initially\\
\rlab{3. Capture forms} & Deploy sign-up forms on the website and blog to gather
names, e-mail addresses and phone numbers\\
\rowcolor{tablealt}
\rlab{4. Define goals} & Determine the mission --- deepen relationships, educate
the audience, or drive sales cycles\\
\rlab{5. Set auto-responders} & Build a sequence of at least six educational,
non-salesy e-mails released automatically to new subscribers\\
\rowcolor{tablealt}
\rlab{6. Add triggers} & Automate new sales cycles based on specific user
actions --- for example clicking a particular link, or abandoning a cart\\
\rlab{7. Monitor metrics} & Review reporting monthly to test subject lines and
improve click-through rate --- and feed the insight back into Node~1\\
\end{mmlongtable}

\section{Worked Plan --- a New Online Business Targeting College Students}

\begin{exambox}[Official Scheme Answer --- 10 Marks]
\pyq{CIE-II, 26 May 2026 --- 10 M}
\textbf{Objective:} increase awareness, engagement and sales among college
students. \textbf{Target audience:} college students aged 18--25.

\begin{enumerate}
\item \sfb{Build the e-mail list} --- collect e-mails through website sign-ups;
      offer discounts and free resources.
\item \sfb{Create attractive content} --- use youthful, creative designs;
      include offers, contests and product updates.
\item \sfb{Segmentation} --- segment users based on interests and purchase
      behaviour.
\item \sfb{Personalisation} --- use customer names and personalised
      recommendations.
\item \sfb{Campaign schedule} --- send weekly newsletters and promotional
      e-mails.
\item \sfb{Mobile optimisation} --- ensure e-mails are mobile-friendly.
\item \sfb{Performance tracking} --- track open rate, click-through rate\ix{click-through rate},
      conversion rate and unsubscribe rate.
\end{enumerate}
\end{exambox}

\section{Segmentation and Personalisation, Step by Step}

\begin{keybox}[High Yield]
``Describe the importance of audience segmentation and personalization in email
marketing. Provide a step-by-step guide on how to segment an audience effectively
and personalize email content based on segmentation criteria.''
\pyq{CIE-II, 24 Jul 2024 --- 10 M}
\end{keybox}

\begin{enumerate}
\item \sfb{Collect the right data at capture} --- ask only for what you will use,
      then enrich progressively over time.
\item \sfb{Choose segmentation criteria} --- demographics (age, location,
      gender); firmographics for B2B (industry, company size, role); behaviour
      (pages viewed, products browsed, past purchases, cart abandonment);
      engagement level (active, lapsing, dormant); lifecycle stage (new
      subscriber, first-time buyer, repeat customer, advocate); and preference
      (the topics and frequency the subscriber selected).
\item \sfb{Build the segments} in the e-mail service provider as \emph{dynamic
      rules}, so that membership updates automatically.
\item \sfb{Design content per segment} --- a different subject line, offer,
      imagery and call to action for each.
\item \sfb{Layer personalisation inside the segment} --- merge tags for the name,
      dynamic product blocks based on browsing history, location-specific offers,
      and send-time optimisation.
\item \sfb{Automate lifecycle triggers} --- welcome series, browse abandonment,
      cart abandonment, post-purchase, replenishment, win-back.
\item \sfb{Test} --- A/B test subject lines and offers \emph{within} each
      segment, since what works for one will not work for another.
\item \sfb{Measure per segment}, and prune or re-engage the segments that
      consistently underperform.
\end{enumerate}

\begin{examplebox}[A Fashion Retailer Segmenting by Past Category]
Subscribers who bought running shoes receive a new-arrivals e-mail featuring
running gear with the subject line ``New for your next 10K'', while subscribers
who bought formal wear receive an entirely different creative.

\smallskip
The same send, two segments --- and typically double the click-through rate of a
single generic blast.
\end{examplebox}

\begin{summarybox}
The e-mail plan has seven nodes --- choose software, build the list, deploy
capture forms, define goals, set auto-responders, add behavioural triggers, and
monitor metrics --- and Node~7 loops back to Node~1. Applied to a scenario, the
plan becomes a specific answer: list building, creative, segmentation,
personalisation, schedule, mobile optimisation and tracking. Segmentation and
personalisation run in eight steps: collect the right data, choose criteria,
build dynamic segments, design per-segment content, layer personalisation inside
the segment, automate lifecycle triggers, test within segments, and measure per
segment.
\end{summarybox}

%%%%%%%%%%%%%%%%%%%%%%%%%%%%%%%%%%%%%%%%%%%%%%%%%%%%%%%%%%%%%%%%%%%%%%%%%%%%%%
\chapter{E-mail Marketing Campaign Analysis}
%%%%%%%%%%%%%%%%%%%%%%%%%%%%%%%%%%%%%%%%%%%%%%%%%%%%%%%%%%%%%%%%%%%%%%%%%%%%%%

\syllabusstrip{E-mail Marketing Campaign Analysis.}

\begin{learningoutcomes}
\begin{itemize}
\item What campaign analysis means.
\item The six-stage measurement funnel, with formulas.
\item How to read each KPI, what a poor reading means, and the fix.
\item The challenges of running e-mail campaigns, with paired solutions.
\item Two concrete strategies for raising open and click rates.
\end{itemize}
\end{learningoutcomes}

\section{What Campaign Analysis Means}

\begin{exambox}[Model Answer --- 2 Marks]
\pyq{CIE-II, 26 May 2026 --- 2 M}
Campaign analysis in e-mail marketing refers to evaluating the performance of
e-mail campaigns using metrics such as open rate, click-through rate, conversion
rate, bounce\ix{bounce} rate and unsubscribe rate, in order to determine campaign
effectiveness and improve future campaigns.
\end{exambox}

\section{The Measurement Funnel}

The funnel runs from the top of the list to the bottom line. Each stage narrows,
and each has its own formula and diagnostic meaning.

\begin{figure}[htbp]
\centering
\begin{tikzpicture}[font=\sffamily\footnotesize]
  % funnel bands, narrowing downward
  \fill[ocre!45]  (-4.6,0)     -- (4.6,0)     -- (4.0,-1.0)  -- (-4.0,-1.0)  -- cycle;
  \fill[ocre!36]  (-4.0,-1.0)  -- (4.0,-1.0)  -- (3.4,-2.0)  -- (-3.4,-2.0)  -- cycle;
  \fill[ocre!28]  (-3.4,-2.0)  -- (3.4,-2.0)  -- (2.8,-3.0)  -- (-2.8,-3.0)  -- cycle;
  \fill[ocre!21]  (-2.8,-3.0)  -- (2.8,-3.0)  -- (2.2,-4.0)  -- (-2.2,-4.0)  -- cycle;
  \fill[ocre!14]  (-2.2,-4.0)  -- (2.2,-4.0)  -- (1.6,-5.0)  -- (-1.6,-5.0)  -- cycle;
  \fill[greenok!30] (-1.6,-5.0) -- (1.6,-5.0) -- (1.0,-6.0)  -- (-1.0,-6.0)  -- cycle;

  \node[text=slate,font=\sffamily\bfseries\scriptsize] at (0,-0.5)  {1.\ \ SENT E-MAILS};
  \node[text=slate,font=\sffamily\bfseries\scriptsize] at (0,-1.5)  {2.\ \ DELIVERY (NON-BOUNCE)};
  \node[text=slate,font=\sffamily\bfseries\scriptsize] at (0,-2.5)  {3.\ \ OPEN RATE};
  \node[text=slate,font=\sffamily\bfseries\scriptsize] at (0,-3.5)  {4.\ \ CLICK-THROUGH RATE};
  \node[text=slate,font=\sffamily\bfseries\scriptsize] at (0,-4.5)  {5.\ \ CONVERSION RATE};
  \node[text=greenok,font=\sffamily\bfseries\scriptsize] at (0,-5.5) {6.\ \ REVENUE};

  \node[anchor=west,text=slatelight,font=\sffamily\scriptsize,align=left] at (4.8,-0.5)
    {list growth rate};
  \node[anchor=west,text=slatelight,font=\sffamily\scriptsize,align=left] at (4.8,-1.5)
    {list hygiene, sender reputation};
  \node[anchor=west,text=slatelight,font=\sffamily\scriptsize,align=left] at (4.8,-2.5)
    {subject line, sender name, timing};
  \node[anchor=west,text=slatelight,font=\sffamily\scriptsize,align=left] at (4.8,-3.5)
    {content relevance, CTA clarity};
  \node[anchor=west,text=slatelight,font=\sffamily\scriptsize,align=left] at (4.8,-4.5)
    {offer quality, landing page};
  \node[anchor=west,text=greenok,font=\sffamily\scriptsize,align=left] at (4.8,-5.5)
    {revenue per message};

  \node[anchor=east,rotate=90,text=ocredark,font=\sffamily\bfseries\scriptsize]
    at (-5.0,-3.0) {THE MARKETING CONTROL ROOM};
\end{tikzpicture}
\caption[The e-mail measurement funnel]%
        {The measurement funnel. Each stage has a distinct diagnosis, so a fall
         in revenue can be traced to exactly one narrowing --- which is why the
         funnel must be read as a system rather than as six separate numbers.}
\label{fig:emailfunnel}
\end{figure}

\begin{mmlongtable}{The six-stage measurement funnel}{tab:funnel}%
{@{}P{0.18\textwidth}P{0.42\textwidth}P{0.33\textwidth}@{}}
\rowcolor{ocre}\thd{Stage} & \thd{Definition or formula} & \thd{What it diagnoses}\\
\midrule
\endfirsthead
\rowcolor{ocre}\thd{Stage} & \thd{Definition or formula} & \thd{What it diagnoses}\\
\midrule
\endhead
\rlab{1. Sent e-mails} & The starting list volume. Metric: list growth rate &
Whether the audience asset is growing or decaying\\
\rowcolor{tablealt}
\rlab{2. Delivery (non-bounce rate)} & $100\% - (\text{bounced} \div
\text{total sent})$. Accounts for invalid addresses or full firewalls &
List hygiene and sender reputation\\
\rlab{3. Open rate} & Tracked when images render on the server or links are
clicked. Calculated from the non-bounce pool & The strength of the subject line,
sender name and send timing\\
\rowcolor{tablealt}
\rlab{4. Click-through rate} & The percentage of opens that clicked a link. Used
to track interests and A/B test variations & The relevance of the content and the
clarity of the call to action\\
\rlab{5. Conversion rate} & The percentage of clicks that completed the ultimate
goal --- purchasing, filling a form & The quality of the offer and the landing
page\\
\rowcolor{tablealt}
\rlab{6. Revenue per message} & $\text{total revenue} \div \text{number of sent
e-mails}$ & The single bottom-line measure of the whole programme\\
\end{mmlongtable}

\begin{exambox}[Define Click-Through Rate --- 2 Marks]
\pyq{CIE-II, 26 May 2026 --- 2 M}
Click-through rate is the percentage of recipients who click on links in an
e-mail compared to the total number of \emph{delivered} e-mails.
\end{exambox}

\begin{warnbox}[Two Denominators, Two Names]
The lecture material defines CTR against \emph{opens}; the scheme of valuation
defines it against \emph{delivered} e-mails. The industry distinguishes these as
\textbf{click-to-open rate} and \textbf{click-through rate} respectively.

\smallskip
State the formula you are using and either is defensible. Stating no formula is
what loses the mark.
\end{warnbox}

\section{Interpreting the KPIs}

\begin{keybox}[High Yield]
``Analyse the effectiveness of an E-mail Marketing campaign using KPIs such as
open rate, click-through rate, unsubscribe rate, and conversion rate.''
\pyq{CIE-II, 26 May 2026 --- 10 M}
Anchor the answer in the six-stage funnel above, so that the KPIs appear as a
connected system rather than a list.
\end{keybox}

\begin{mmlongtable}{Reading the e-mail KPIs}{tab:emailkpi}%
{@{}P{0.157\textwidth}P{0.235\textwidth}P{0.255\textwidth}P{0.265\textwidth}@{}}
\rowcolor{ocre}\thd{KPI} & \thd{What it measures} & \thd{What a poor reading means} & \thd{Fix}\\
\midrule
\endfirsthead
\rowcolor{ocre}\thd{KPI} & \thd{What it measures} & \thd{What a poor reading means} & \thd{Fix}\\
\midrule
\endhead
\rlab{Open rate} & How many users opened the e-mail & Weak subject line, wrong
send time, poor sender reputation & Test subject lines; clean the list;
authenticate the domain\\
\rowcolor{tablealt}
\rlab{Click-through rate} & How many users clicked links in the e-mail & Content
not engaging, or the call to action unclear & Improve the offer; make one CTA
dominant; segment better\\
\rlab{Unsubscribe rate} & Users opting out of e-mails & Content is irrelevant, or
frequency is too high & Send relevant content; offer a preference centre instead
of only ``unsubscribe''\\
\rowcolor{tablealt}
\rlab{Conversion rate} & Users completing the desired action & The offer or
landing page is failing, not the e-mail & Fix message match on the landing page;
reduce form friction\\
\rlab{Bounce rate} & Undelivered e-mails & Stale or purchased list; server blocks
& Remove hard bounces immediately; use double opt-in\\
\rowcolor{tablealt}
\rlab{Spam complaint rate} & Recipients marking the e-mail as junk & Consent is
unclear, or expectations were not set & Set expectations at sign-up; make
unsubscribe obvious\\
\end{mmlongtable}

\begin{keybox}[The Diagnostic Logic]
Read the funnel downward and the failure locates itself. Clicks but no
conversions is \emph{not} an e-mail problem --- the e-mail did its job and handed
the visitor to a page that failed. Opens but no clicks \emph{is} an e-mail
problem: the subject line promised something the body did not deliver.
\end{keybox}

\section{Challenges and Solutions}

\begin{exambox}[Official Scheme Answer --- 10 Marks]
\pyq{CIE-II, 26 May 2026 --- 10 M}
\end{exambox}

\begin{mmlongtable}{E-mail marketing challenges and solutions}{tab:challenges}%
{@{}P{0.38\textwidth}P{0.55\textwidth}@{}}
\rowcolor{ocre}\thd{Challenge} & \thd{Solution}\\
\midrule
\endfirsthead
\rowcolor{ocre}\thd{Challenge} & \thd{Solution}\\
\midrule
\endhead
\rlab{Low open rates} & Use attractive subject lines\\
\rowcolor{tablealt}
\rlab{Spam filtering issues} & Avoid spam words and maintain sender reputation\\
\rlab{High unsubscribe rates} & Send relevant and useful content\\
\rowcolor{tablealt}
\rlab{Poor e-mail design} & Use responsive e-mail templates\\
\rlab{Lack of personalisation} & Segment the audience and personalise e-mails\\
\rowcolor{tablealt}
\rlab{Inactive subscribers} & Run re-engagement campaigns\\
\end{mmlongtable}

\noindent Additional challenges worth naming: deliverability and inbox
placement; \keytermas{list decay}{list decay} of roughly 20--25\% per year;
rendering inconsistency across clients; privacy changes such as mail privacy
protection that inflate open rates and make them unreliable; and a rising
regulatory burden.

\section{Raising Open and Click Rates}

\begin{exambox}[Two Strategies --- 2 Marks, for a Fashion Brand's Summer Sale]
\pyq{CIE-II Quiz, 24 Jul 2024 --- 2 M}
\begin{itemize}
\item \sfb{To increase open rate:} write a specific, urgency-bearing,
      personalised subject line and A/B test it --- for example ``Priya, your
      summer edit is 40\% off --- today only'' --- and optimise the send time to
      the segment's historical open window.
\item \sfb{To increase click-through rate:} use a single dominant, benefit-led
      call-to-action button above the fold, supported by a strong hero product
      image and segmented product recommendations based on the recipient's
      browsing history.
\end{itemize}
\end{exambox}

\begin{summarybox}
Campaign analysis evaluates e-mail performance on open rate, click-through rate,
conversion rate, bounce rate and unsubscribe rate in order to improve future
campaigns. The measurement funnel has six stages --- sent, delivered, opened,
clicked, converted, revenue --- each with a formula and a distinct diagnosis, so
a fall in revenue traces to exactly one narrowing. Each KPI has a characteristic
failure and a matching fix, and the key diagnostic distinction is that clicks
without conversions indicts the landing page while opens without clicks indicts
the e-mail. Six standard challenges pair with six solutions, and to those should
be added deliverability, list decay of 20--25\% a year, rendering inconsistency
and privacy changes that make open rate unreliable.
\end{summarybox}

%%%%%%%%%%%%%%%%%%%%%%%%%%%%%%%%%%%%%%%%%%%%%%%%%%%%%%%%%%%%%%%%%%%%%%%%%%%%%%
\chapter{Measuring Conversions and Keeping Up}
%%%%%%%%%%%%%%%%%%%%%%%%%%%%%%%%%%%%%%%%%%%%%%%%%%%%%%%%%%%%%%%%%%%%%%%%%%%%%%

\syllabusstrip{Measuring Conversions \& keeping up.}

\begin{learningoutcomes}
\begin{itemize}
\item How to track conversions that happen offline.
\item How to integrate web analytics with e-mail marketing.
\item The E-mail 2.0 technology stack.
\item Responsive e-mail design --- the case, the method and the critical view.
\item The four pillars of the e-mail engine.
\end{itemize}
\end{learningoutcomes}

\section{Bridging the Gap --- Offline Tracking}

When open rates and click-through rates are not enough to track overall marketing
success, e-mail requires physical participation mechanisms.

\begin{mmlongtable}{Tracking conversions that happen offline}{tab:offline}%
{@{}P{0.22\textwidth}P{0.71\textwidth}@{}}
\rowcolor{ocre}\thd{Mechanism} & \thd{How it is tracked}\\
\midrule
\endfirsthead
\rowcolor{ocre}\thd{Mechanism} & \thd{How it is tracked}\\
\midrule
\endhead
\rlab{In-store purchases} & By requesting the customer's e-mail address at the
physical point of sale and matching it against the campaign list\\
\rowcolor{tablealt}
\rlab{Phone calls} & By providing dedicated phone numbers exclusive to the
e-mail campaign\\
\rlab{Event attendance} & By issuing special entrance codes that can only be
accessed and viewed from within the e-mail itself\\
\end{mmlongtable}

\begin{keybox}[The Digital Equivalents of the Same Principle]
Unique promotion codes per campaign; UTM-tagged links feeding web analytics; and
printable or scannable QR coupons that carry a campaign identifier. In every case
the mechanism is the same: give this campaign, and only this campaign, a token
that can be counted later.
\end{keybox}

\section{Integrating Web Analytics and E-mail Marketing}

\begin{exambox}[Official Scheme Answer --- 10 Marks]
\pyq{CIE-II, 26 May 2026 --- 10 M}
\begin{enumerate}
\item \sfb{Collect customer data} --- use web analytics tools to track user
      behaviour, page visits and purchase history.
\item \sfb{Audience segmentation} --- segment users based on interests and
      browsing behaviour.
\item \sfb{Personalised e-mail campaigns} --- send customised e-mails with
      relevant product recommendations.
\item \sfb{Track campaign performance} --- use open rate, click-through rate,
      conversion rate and bounce rate.
\item \sfb{Retargeting} --- send follow-up e-mails to abandoned-cart users.
\item \sfb{Continuous optimisation} --- analyse campaign data and improve content
      and timing.
\end{enumerate}

\smallskip
\textbf{Benefits:} better customer engagement; higher conversion rates; improved
customer retention\ix{customer retention}; increased sales and revenue.
\end{exambox}

\section{The E-mail 2.0 Technology Stack}

\begin{mmlongtable}{The e-mail technology stack}{tab:emailstack}%
{@{}P{0.24\textwidth}P{0.69\textwidth}@{}}
\rowcolor{ocre}\thd{Category} & \thd{Tools}\\
\midrule
\endfirsthead
\rowcolor{ocre}\thd{Category} & \thd{Tools}\\
\midrule
\endhead
\rlab{CRM and all-in-one automation} & HubSpot (automated workflows),
Infusionsoft (purchase behaviour and task management), ActiveCampaign (lifecycle
automation), Contactually (social and inbox vetting)\\
\rowcolor{tablealt}
\rlab{Drip campaigns and nurturing} & Drip (lightweight workflow builder),
ConvertKit (automation rules for publishers), Emma (personalisation via CRM
sync), Campaign Monitor (interactive analytics), Sendloop\\
\rlab{Creation and standard sending} & MailChimp (abandoned carts, groups),
GetResponse (calendar-integrated auto-responders), AWeber (new-subscriber
catch-up), Reach Mail (simultaneous comparison of five campaigns)\\
\rowcolor{tablealt}
\rlab{Specialised utility} & Litmus (one-click render testing across around 40
clients), BombBomb (video-powered e-mail), Gumroad (digital product updates),
Intercom (in-app conversations)\\
\rlab{Sales tracking plugins} & GetNotify (invisible open tracking), ToutApp
(playbooks and CRM sync), Yesware (prescriptive analytics, team tracking),
Clearslide (web engagement tracking)\\
\end{mmlongtable}

\section{Responsive E-mail Design}

\begin{keybox}[High Yield]
``Evaluate the effectiveness of responsive email design in maintaining a
competitive edge in digital marketing.''
\pyq{SEE, Jun/Jul 2025, Q6a --- 8 M}
The verb is \emph{evaluate}, so the answer must reach a judgement rather than
merely describe the technique.
\end{keybox}

\subsection{The case for it}

The majority of e-mail opens now occur on mobile, and a non-responsive e-mail
requiring pinch-and-zoom is typically deleted within seconds. Responsive design\ix{responsive design}
therefore protects open-to-click conversion, reduces unsubscribes, improves
accessibility, and preserves brand perception. It also improves deliverability
indirectly, because engagement signals feed inbox placement algorithms.

\subsection{How it is achieved}

\begin{itemize}
\item Fluid single-column layouts and media queries.
\item Scalable images.
\item Body text of at least 14 to 16 pixels.
\item Tap targets of at least 44 by 44 pixels.
\item ``Bulletproof'' call-to-action buttons built in HTML rather than as images.
\item Alt text, for image-blocked clients.
\end{itemize}

\subsection{The critical view}

\begin{warnbox}[Responsiveness Is Table Stakes, Not a Differentiator]
Every competent competitor has responsive e-mail, so it prevents disadvantage
more than it creates advantage.

\smallskip
Genuine competitive edge comes from \textbf{segmentation, personalisation,
behavioural triggers and offer quality}. Responsive design is the
\emph{precondition} that allows those to work.

\smallskip
It also adds testing overhead across dozens of clients --- hence render-testing
tools --- and dark-mode rendering remains inconsistent.
\end{warnbox}

\section{Operating the Engine --- Four Pillars}

\begin{figure}[htbp]
\centering
% The pillar shafts are drawn as fixed rectangles and the text is placed as
% separately anchored nodes inside them, so no amount of text can push a shaft
% through the capital band above it.
\begin{tikzpicture}[font=\sffamily\footnotesize,
    hd/.style={anchor=north,align=center,text width=2.40cm,
               font=\sffamily\bfseries\scriptsize,text=ocredark},
    sb/.style={anchor=north,align=center,text width=2.40cm,
               font=\sffamily\scriptsize\itshape,text=slatelight},
    bd/.style={anchor=north,align=center,text width=2.40cm,
               font=\sffamily\scriptsize,text=slate}]

  % capital and base
  \fill[ocre]  (-0.35,4.80) rectangle (11.50,5.10);
  \fill[slate] (-0.35,-0.40) rectangle (11.50,-0.08);
  \node[anchor=north,text=slate,font=\sffamily\bfseries\scriptsize]
    at (5.58,-0.55) {A TARGETED, MEASURABLE, HIGH-ROI ECOSYSTEM};

  % ---- pillar 1
  \fill[ocrepale] (0.0,0) rectangle (2.6,4.75);
  \draw[ocre,line width=1pt] (0.0,0) rectangle (2.6,4.75);
  \node[hd] at (1.3,4.62) {PILLAR 1\\ THE FUEL};
  \node[sb] at (1.3,3.85) {the opt-in list};
  \node[bd] at (1.3,3.42) {micro-focused, permission-based audiences that
     eliminate legal risk and guarantee engagement};

  % ---- pillar 2
  \fill[ocrepale] (2.85,0) rectangle (5.45,4.75);
  \draw[ocre,line width=1pt] (2.85,0) rectangle (5.45,4.75);
  \node[hd] at (4.15,4.62) {PILLAR 2\\ THE PAYLOAD};
  \node[sb] at (4.15,3.85) {content arsenal};
  \node[bd] at (4.15,3.42) {deploy the exact right template: promotional,
     educational, survey};

  % ---- pillar 3
  \fill[ocrepale] (5.70,0) rectangle (8.30,4.75);
  \draw[ocre,line width=1pt] (5.70,0) rectangle (8.30,4.75);
  \node[hd] at (7.0,4.62) {PILLAR 3\\ THE ENGINE};
  \node[sb] at (7.0,3.85) {triggers and tech};
  \node[bd] at (7.0,3.42) {move from static blasts to dynamic behavioural
     lifecycle automation};

  % ---- pillar 4
  \fill[ocrepale] (8.55,0) rectangle (11.15,4.75);
  \draw[ocre,line width=1pt] (8.55,0) rectangle (11.15,4.75);
  \node[hd] at (9.85,4.62) {PILLAR 4\\ THE DASHBOARD};
  \node[sb] at (9.85,3.85) {the metrics};
  \node[bd] at (9.85,3.42) {obsess over the funnel, from non-bounce delivery to
     offline revenue};
\end{tikzpicture}
\caption[The four pillars of the e-mail engine]%
        {The four pillars. Remove any one and the engine stalls: a perfect
         template sent to a purchased list is illegal, and a perfectly automated
         programme without a dashboard cannot be improved.}
\label{fig:fourpillars}
\end{figure}

\begin{mmlongtable}{The four pillars}{tab:pillars}%
{@{}P{0.28\textwidth}P{0.65\textwidth}@{}}
\rowcolor{ocre}\thd{Pillar} & \thd{What it means}\\
\midrule
\endfirsthead
\rowcolor{ocre}\thd{Pillar} & \thd{What it means}\\
\midrule
\endhead
\rlab{1 --- The Fuel (opt-in list)} & Cultivate micro-focused, permission-based
audiences, to eliminate the legal risk and guarantee engagement\\
\rowcolor{tablealt}
\rlab{2 --- The Payload (content arsenal)} & Deploy the exact right template ---
promotional to drive revenue, educational to build trust, surveys to gather
intelligence\\
\rlab{3 --- The Engine (triggers and tech)} & Use the technology stack to move
from static blasts to dynamic, behavioural, lifecycle automation\\
\rowcolor{tablealt}
\rlab{4 --- The Dashboard (metrics)} & Obsess over the measurement funnel ---
from non-bounce delivery down to total offline revenue generation\\
\end{mmlongtable}

\begin{keybox}[The Closing Idea]
E-mail marketing is no longer a megaphone. It is a targeted, measurable, high-ROI
ecosystem.
\end{keybox}

\begin{summarybox}
Conversions that happen offline are tracked by giving each campaign a countable
token: an e-mail address captured at the point of sale, a dedicated phone number,
an e-mail-only entrance code, or their digital equivalents in promotion codes,
UTM tags and QR coupons. Web analytics and e-mail integrate in six steps ---
collect behavioural data, segment, personalise, track performance, retarget
abandoners, and optimise continuously --- yielding better engagement, higher
conversion, improved retention and increased revenue. The technology stack spans
CRM automation, drip nurturing, standard sending, specialised utilities and sales
tracking. Responsive design protects open-to-click conversion and is now a
precondition rather than a differentiator; the real edge lies in segmentation,
personalisation, triggers and offer quality. Four pillars operate the engine:
the opt-in list, the content arsenal, the trigger technology and the metrics
dashboard.
\end{summarybox}

%%%%%%%%%%%%%%%%%%%%%%%%%%%%%%%%%%%%%%%%%%%%%%%%%%%%%%%%%%%%%%%%%%%%%%%%%%%%%%
%  Unit III appendices
%%%%%%%%%%%%%%%%%%%%%%%%%%%%%%%%%%%%%%%%%%%%%%%%%%%%%%%%%%%%%%%%%%%%%%%%%%%%%%

\startappendices

\chapter{Syllabus-to-Chapter Concordance}

This appendix exists so that coverage can be \emph{audited} rather than assumed.
The left-hand column reproduces the Unit~III syllabus descriptor phrase by
phrase, in the order the official document prints it. The right-hand column gives
the chapter and section that carries it. Every phrase resolves to a location, and
no numbered chapter of this volume contains material that is not traceable to a
phrase in this table.

\begin{mmlongtable}{Concordance for Unit III}{tab:concordance3}%
{@{}P{0.36\textwidth}P{0.57\textwidth}@{}}
\rowcolor{ocre}\thd{Syllabus phrase} & \thd{Where it is covered}\\
\midrule
\endfirsthead
\rowcolor{ocre}\thd{Syllabus phrase} & \thd{Where it is covered}\\
\midrule
\endhead
\multicolumn{2}{@{}l}{\bfseries\color{ocredark} First half --- Web Analytics \& Google Analytics}\\
\midrule
\rlab{Web Analytics} \newline {\scriptsize the heading concept} & Chapter~1
entire. Definition and importance \S\,1.1; traditional against web-analytics
based analysis \S\,1.2; the three metric families \S\,1.3; core metrics
\S\,1.4; bounce rate \S\,1.5; the dual lens \S\,1.6; the role of web analytics
\S\,1.7; benefits \S\,1.8; KPI alignment \S\,1.9; A/B testing and CRO
\S\,1.10.\\
\rowcolor{tablealt}
\rlab{Google Analytics Tools} & Chapter~2 entire. Purpose \S\,2.1; key features
\S\,2.2; acquisition channels \S\,2.3; report families \S\,2.4; setting up and
tracking goals \S\,2.5; Google Tag Manager \S\,2.6.\\
\rlab{Web Analytics Tools} & Chapter~3 entire. Tools other than Google Analytics
\S\,3.1; combining quantitative and qualitative tools \S\,3.2.\\
\midrule
\multicolumn{2}{@{}l}{\bfseries\color{ocredark} Second half --- E-mail Marketing}\\
\midrule
\rlab{Effective E-mail Campaigns} & Chapter~4 entire. Definition \S\,4.1; E-mail
1.0 against 2.0 \S\,4.2; advantages \S\,4.3; disadvantages \S\,4.4; opt-in
against spam \S\,4.5; anatomy of an effective e-mail \S\,4.6; the nine-type
content arsenal \S\,4.7; transactional against promotional \S\,4.8.\\
\rowcolor{tablealt}
\rlab{E-mail Plan} & Chapter~5 entire. The seven nodes \S\,5.1; a worked plan
\S\,5.2; segmentation and personalisation \S\,5.3.\\
\rlab{E-mail Marketing Campaign Analysis} & Chapter~6 entire. What campaign
analysis means \S\,6.1; the measurement funnel \S\,6.2; interpreting the KPIs
\S\,6.3; challenges and solutions \S\,6.4; raising open and click rates
\S\,6.5.\\
\rowcolor{tablealt}
\rlab{Measuring Conversions \& keeping up} & Chapter~7 entire. Offline tracking
\S\,7.1; integrating analytics with e-mail \S\,7.2; the technology stack
\S\,7.3; responsive design \S\,7.4; the four pillars \S\,7.5.\\
\end{mmlongtable}

\begin{keybox}[No Off-Syllabus Appendix in This Volume]
Unlike Volumes~1 and~2, every topic in the source material for Unit~III maps to a
phrase in the unit descriptor. Nothing had to be set aside, so this volume has no
appendix of supplementary examined topics.
\end{keybox}

\chapter{Previous-Year Question Bank --- Unit III}

Every question in this appendix is reproduced from an actual RVCE Marketing
Management Continuous Internal Evaluation or Semester End Examination paper for
this course. Answers and pointers follow the official schemes of valuation where
those were released.

\section*{Part A --- Two-Mark Questions}
\addcontentsline{toc}{section}{Part A --- Two-Mark Questions}

\begin{enumerate}[leftmargin=2.2em,itemsep=1.1ex]

\item \textbf{Define Web Analytics and mention its importance in digital
      marketing activities.} \pyqr{CIE-II, 26 May 2026 --- 2 M}\\
      Web analytics is the process of collecting, measuring, analysing and
      reporting website data to understand and optimise web usage. Importance: it
      helps understand customer behaviour, measures campaign effectiveness,
      improves website performance and supports better marketing decisions.
      \hfill\emph{See} \S\,1.1.

\item \textbf{What is Bounce\ix{bounce} Rate in Web Analytics and how does it affect website
      performance evaluation? / How bounce rate affects website performance
      assessment.} \pyqr{CIE-II, 26 May 2026; SEE, Jun/Jul 2025 --- 2 M}\\
      Bounce rate is the percentage of visitors who leave a website after viewing
      only one page without any interaction. A high bounce rate signals poor
      content relevance, slow loading, confusing design or mismatched traffic,
      and correlates with lower conversions --- so it acts as an early diagnostic
      of both content quality and traffic quality. \hfill\emph{See} \S\,1.5.

\item \textbf{List any four important features available in Google Analytics\ix{Google Analytics} for
      website performance tracking.} \pyqr{CIE-II, 26 May 2026 --- 2 M}\\
      Traffic source analysis; bounce rate measurement; user demographics
      tracking; conversion and goal tracking. \hfill\emph{See} \S\,2.2.

\item \textbf{What is the main purpose of using Google Analytics?}
      \pyqr{CIE-II, 7 May 2025 --- 2 M}\\
      To collect, measure and report data about website users --- who they are,
      how they arrived, what they did and whether they converted --- so that
      businesses can understand behaviour, evaluate marketing performance and
      make evidence-based decisions. \hfill\emph{See} \S\,2.1.

\item \textbf{Name any two popular web analytics tools other than Google
      Analytics.} \pyqr{CIE-II, 7 May 2025 --- 2 M}\\
      Adobe Analytics and Matomo. Also acceptable: Webtrends, Mixpanel, Hotjar\ix{Hotjar},
      Microsoft Clarity, Kissmetrics, Amplitude, SimilarWeb.
      \hfill\emph{See} \S\,3.1.

\item \textbf{Define ``Dimensions'' and ``Metrics'' in Google Analytics.}
      \pyqr{SEE, Oct/Nov 2023 --- 2 M}\\
      Dimensions are attributes of the data --- qualitative descriptors such as
      city, device, browser, source and medium, or page path; they form the rows
      of a report. Metrics are quantitative measurements --- sessions, users,
      pageviews, bounce rate, revenue; they form the columns.
      \hfill\emph{See} \S\,1.4.

\item \textbf{What is meant by campaign analysis in E-mail Marketing?}
      \pyqr{CIE-II, 26 May 2026 --- 2 M}\\
      Evaluating the performance of e-mail campaigns using metrics such as open
      rate, click-through rate\ix{click-through rate}, conversion rate, bounce rate and unsubscribe
      rate, to determine effectiveness and improve future campaigns.
      \hfill\emph{See} \S\,6.1.

\item \textbf{Define Click Through Rate (CTR) in E-mail Marketing.}
      \pyqr{CIE-II, 26 May 2026 --- 2 M}\\
      CTR is the percentage of recipients who click on links in an e-mail
      compared to the total number of delivered e-mails.
      \hfill\emph{See} \S\,6.2.

\item \textbf{List any four advantages of E-mail Marketing over traditional
      marketing methods.} \pyqr{CIE-II, 26 May 2026 --- 2 M}\\
      Low cost; global reach; easy performance tracking; personalised
      communication. \hfill\emph{See} \S\,4.3.

\item \textbf{Why is segmentation important in e-mail marketing campaigns?}
      \pyqr{SEE, Jun/Jul 2025 --- 2 M}\\
      Segmentation divides the list by demographics, behaviour, engagement or
      lifecycle stage so that each group receives relevant content. It raises open
      and click rates, reduces unsubscribes and spam complaints, protects
      deliverability and enables the personalisation that drives conversion.
      \hfill\emph{See} \S\,4.5.

\item \textbf{List any two key components of a successful email marketing
      campaign.} \pyqr{SEE, Oct/Nov 2023 --- 2 M}\\
      A permission-based, well-segmented opt-in list, and a compelling subject
      line with a single clear call to action\ix{call to action}. Also acceptable: personalised
      relevant content; mobile-responsive design; measurable KPIs.
      \hfill\emph{See} \S\,4.6.

\item \textbf{After conducting an email marketing campaign for a travel agency,
      outline two specific metrics you would analyze to assess the campaign's
      performance.} \pyqr{CIE-II Quiz, 24 Jul 2024 --- 2 M}\\
      Click-through rate, to judge whether the destinations and offers presented
      were relevant enough to earn a click; and conversion rate --- bookings or
      enquiry-form completions per click --- to judge whether the campaign
      produced business value rather than merely interest.
      \hfill\emph{See} \S\,6.3.

\item \textbf{You are tasked with creating an effective email campaign to promote
      a summer sale for a fashion brand. Describe two strategies you would employ
      to increase open rates and click-through rates.}
      \pyqr{CIE-II Quiz, 24 Jul 2024 --- 2 M}\\
      A personalised, urgency-bearing, A/B-tested subject line with an optimised
      send time to raise opens; and a single dominant benefit-led call to action
      above the fold with segmented product recommendations to raise clicks.
      \hfill\emph{See} \S\,6.5.

\end{enumerate}

\section*{Part B --- Eight- and Ten-Mark Questions}
\addcontentsline{toc}{section}{Part B --- Eight- and Ten-Mark Questions}

\begin{enumerate}[leftmargin=2.2em,itemsep=1.4ex]

\item \textbf{Differentiate between traditional marketing analysis and web
      analytics-based marketing analysis with suitable examples.}
      \pyqr{CIE-II, 26 May 2026 --- 10 M}\\
      Use the comparison table --- data collection, speed, accuracy, customer
      tracking, cost, measurement, sample, attribution and feedback loop --- each
      with an example. \hfill\emph{See} \S\,1.2.

\item \textbf{Explain the role of web analytics in digital marketing. How can
      businesses benefit from analysing website data?}
      \pyqr{CIE-II, 7 May 2025 --- 10 M}\\
      The eight roles and the seven benefits. \hfill\emph{See} \S\,1.7 and
      \S\,1.8.

\item \textbf{Develop a basic strategy to set up Google Analytics for a new
      e-commerce website. What key metrics should be tracked?}
      \pyqr{CIE-II, 7 May 2025 --- 10 M}\\
      \emph{Setup strategy:} (1) create the property and data stream; (2) deploy
      the tag through Google Tag Manager\ix{Google Tag Manager} rather than hard-coding; (3) define the
      measurement plan --- what business questions must the data answer; (4)
      configure enhanced e-commerce events (view item, add to cart, begin
      checkout, purchase); (5) mark purchase and lead events as conversions and
      assign values; (6) build the checkout funnel exploration; (7) apply UTM
      tagging conventions to every campaign; (8) link the advertising platform
      and Search Console; (9) configure internal-traffic and bot filters; (10)
      set up audiences for remarketing; (11) build a dashboard; (12) validate in
      debug view and real time before trusting the data.
      \emph{Key metrics:} sessions and users by channel; new against returning;
      bounce or engagement rate; average session duration; product views,
      add-to-cart rate and cart-abandonment rate; checkout completion rate;
      conversion rate; average order value; revenue; return on advertising spend;
      and cost per acquisition by channel.
      \hfill\emph{See} \S\,2.2 and \S\,2.5.

\item \textbf{How can businesses align Web Analytics KPIs with their
      organizational goals? Illustrate with examples.}
      \pyqr{CIE-I, 16 Apr 2026 --- 10 M}\\
      The five-step method with the goal-to-KPI mapping and a SMART\ix{SMART} example.
      \hfill\emph{See} \S\,1.9.

\item \textbf{Discuss the concept of ``Acquisition channels'' in Google
      Analytics. How can business use this information to optimize their
      marketing efforts?} \pyqr{SEE, Oct/Nov 2023, Q5a --- 8 M}\\
      Define acquisition channels, list them with meaning and optimisation
      action, and explain the ways businesses act on the data --- above all,
      judging a channel by revenue rather than by traffic.
      \hfill\emph{See} \S\,2.3.

\item \textbf{Explain the significance of A/B testing\ix{A/B testing} and conversion rate
      optimization in Web Analytics. How can these practices help improve website
      performance?} \pyqr{SEE, Oct/Nov 2023, Q5b --- 8 M}\\
      Definitions, significance, what to test, the six-step process and the
      tools. \hfill\emph{See} \S\,1.10.

\item \textbf{List and explain the various reports that can be generated using
      Google Analytics.} \pyqr{SEE, Oct/Nov 2023, Q6b --- 8 M}\\
      Real-time, audience and user, acquisition, behaviour and engagement,
      conversions and monetisation, attribution, and custom reports and
      explorations. \hfill\emph{See} \S\,2.4.

\item \textbf{Describe the process of setting up and tracking goals in Google
      Analytics. How can businesses use goal tracking to measure website
      performance?} \pyqr{SEE, Oct/Nov 2023, Q10b --- 6 M}\\
      The eight-step process, and the point that goals convert traffic reports
      into ROI reports. \hfill\emph{See} \S\,2.5.

\item \textbf{Design a simple marketing plan for an e-learning platform using
      insights gathered from web analytics tools.}
      \pyqr{SEE, Jun/Jul 2025, Q5a --- 8 M}\\
      \emph{Insight gathering:} analytics shows that most traffic arrives from
      organic search\ix{organic search} on mobile, that the highest-engagement pages are free sample
      lessons, and that 70\% of drop-off occurs on the payment page; session
      recordings show users hesitating at a nine-field sign-up form.
      \emph{Plan:} (1) objective --- increase paid enrolments by 25\% in one
      quarter; (2) audience --- segment into browsers, free-trial users and
      lapsed users; (3) channels --- SEO around ``learn X online'' long-tail
      queries, video sample lessons, and a re-engagement e-mail sequence; (4)
      content --- the free sample lesson as the lead magnet, since analytics
      identified it as the strongest engagement asset; (5) conversion fixes ---
      cut the sign-up form to four fields, add trust signals and a money-back
      guarantee at the payment step, and mobile-optimise checkout; (6) e-mail ---
      trigger a nurture sequence on free-trial start and an abandoned-checkout
      reminder; (7) measurement --- track enrolment conversion rate, cost per
      enrolment, checkout completion rate and course-completion retention;
      review fortnightly. \hfill\emph{See} Chapter~1 and Chapter~5.

\item \textbf{Analyze a failed e-mail marketing campaign and identify areas where
      marketing research could have improved outcomes.}
      \pyqr{SEE, Jun/Jul 2025, Q5b --- 8 M}\\
      Diagnose against the funnel. \emph{High bounce rate} \ra{} the list was
      purchased or stale; research into list source and hygiene would have
      prevented it. \emph{Low open rate} \ra{} the subject line and sender name
      were untested; subject-line A/B research would have identified a winner.
      \emph{Low CTR despite opens} \ra{} content was irrelevant to the segment;
      audience research and segmentation would have matched offer to interest.
      \emph{Clicks but no conversions} \ra{} the landing page did not match the
      e-mail's promise; message-match testing and usability\ix{usability} research would have
      caught it. \emph{High unsubscribes and spam complaints} \ra{} expectations
      were never set at sign-up and frequency was too high; preference research
      would have fixed both.

      \smallskip
      \emph{Conclusion:} the failure was not a creative failure but a
      \textbf{research failure} --- the campaign was designed from assumption
      rather than from audience data. \hfill\emph{See} \S\,6.2 and \S\,6.3.

\item \textbf{Describe the importance of audience segmentation and
      personalization in email marketing. Provide a step-by-step guide on how to
      segment an audience effectively and personalize email content based on
      segmentation criteria.} \pyqr{CIE-II, 24 Jul 2024 --- 10 M}\\
      The eight-step guide with the fashion-retailer example, preceded by the
      importance argument. \hfill\emph{See} \S\,5.3 and \S\,4.5.

\item \textbf{Imagine you are a marketing manager for a startup selling organic
      skincare products. Develop a case study outlining the steps you would take
      to create and execute an effective e-mail campaign to launch a new product
      line. Include specific goals, target audience segmentation, content
      strategy, and metrics for success.}
      \pyqr{CIE-II, 24 Jul 2024 --- 10 M}\\
      \emph{Goals:} 3,000 pre-launch sign-ups; 25\% open rate; 4\% conversion on
      the launch send; \rupee 8 lakh revenue in the launch fortnight.
      \emph{Segmentation:} existing customers by past category (cleansers,
      serums, sun care); engaged non-buyers; customers lapsed over 180 days; and
      new sign-ups from the pre-launch teaser.
      \emph{Content strategy --- a five-e-mail sequence:} (1) a teaser to the
      full list, building anticipation with an ingredient story; (2) an
      educational e-mail on the skin concern the range solves, establishing
      authority; (3) the launch announcement with a hero product image and a
      single call to action, personalised by segment; (4) a social-proof e-mail
      carrying testimonials and before-and-after images; (5) an urgency e-mail
      for non-openers and cart abandoners with a 48-hour launch offer. Each is
      mobile-responsive, single-topic, with one dominant call to action.
      \emph{Execution:} double opt-in capture, dynamic product blocks by segment,
      send-time optimisation, and A/B tests on subject lines for e-mails 3 and 5.
      \emph{Metrics:} list growth rate, non-bounce rate, open rate,
      click-through rate, conversion rate, revenue per message, unsubscribe and
      spam-complaint rate --- reviewed after each send so the sequence is
      optimised mid-flight. \hfill\emph{See} Chapter~5 and Chapter~6.

\item \textbf{Evaluate the challenges faced in maintaining successful E-mail
      Marketing campaigns and suggest suitable solutions to improve campaign
      performance.} \pyqr{CIE-II, 26 May 2026 --- 10 M}\\
      The six challenges with paired solutions, plus deliverability, list decay,
      rendering inconsistency and privacy changes. \hfill\emph{See} \S\,6.4.

\item \textbf{Create a strategy to integrate Web Analytics and E-mail Marketing
      for improving customer engagement and business growth.}
      \pyqr{CIE-II, 26 May 2026 --- 10 M}\\
      The six-step strategy and the four benefits. \hfill\emph{See} \S\,7.2.

\item \textbf{Analyse the effectiveness of an E-mail Marketing campaign using
      KPIs such as open rate, click-through rate, unsubscribe rate, and
      conversion rate.} \pyqr{CIE-II, 26 May 2026 --- 10 M}\\
      Define each KPI, state what a high or low reading indicates, and give the
      corrective action. Anchor the answer in the six-stage measurement funnel so
      that the KPIs are shown as a connected system rather than a list.
      \hfill\emph{See} \S\,6.2 and \S\,6.3.

\item \textbf{Design an E-mail Marketing plan for a newly launched online business
      targeting college students.} \pyqr{CIE-II, 26 May 2026 --- 10 M}\\
      The official seven-point scheme answer. \hfill\emph{See} \S\,5.2.

\item \textbf{Evaluate the effectiveness of responsive email design in
      maintaining a competitive edge in digital marketing.}
      \pyqr{SEE, Jun/Jul 2025, Q6a --- 8 M}\\
      The case for it, how it is achieved, and the critical view that it is now
      table stakes rather than a differentiator. \hfill\emph{See} \S\,7.4.

\item \textbf{Critically evaluate the contribution of Google Tag Manager in
      optimizing marketing channels for global companies.}
      \pyqr{SEE, Jun/Jul 2025, Q6b --- 8 M}\\
      The six contributions and the four critical limitations.
      \hfill\emph{See} \S\,2.6.

\item \textbf{Demonstrate how tools like Google Analytics and Hotjar can be used
      to develop insights into user behavior and improve marketing decisions.}
      \pyqr{SEE, Jun/Jul 2025, Q7a --- 8 M}\\
      Analytics gives the \emph{what} --- which channels deliver traffic, which
      landing pages bounce, where users drop out of checkout, which segments
      convert, and what each channel costs per acquisition. Behavioural tools
      give the \emph{why} --- heatmaps show which elements attract attention,
      scroll maps show whether users ever reach the call to action, click maps
      expose rage clicks on non-clickable elements, session recordings show
      hesitation in real time, and on-page surveys capture the reason for
      abandonment in the user's own words. Then give the combined workflow and
      the decisions it improves. \hfill\emph{See} \S\,3.2.

\end{enumerate}

\section*{How Unit III Is Examined}
\addcontentsline{toc}{section}{How Unit III Is Examined}

\begin{keybox}[The Pattern Across Papers]
In the Semester End Examination, Unit~III is Questions~5 and~6 with internal
choice --- 16 marks in two eight-mark sub-divisions --- plus two or three Part~A
objective questions. The second Continuous Internal Evaluation test draws heavily
on it.

\smallskip
The highest-probability topics are the \textbf{e-mail measurement funnel and KPI
analysis}, \textbf{designing an e-mail marketing plan} for a given scenario,
\textbf{segmentation and personalisation}, \textbf{Google Analytics reports, goals
and acquisition channels}, and \textbf{traditional against web-analytics-based
analysis}.

\smallskip
Note that almost every Unit~III question is \emph{applied} --- you will be given a
business and asked to build or diagnose something --- so always answer with a
named framework (the seven-node plan, the six-stage funnel, the eight-step goal
setup) rather than general commentary.
\end{keybox}

\chapter{Glossary of Key Terms}

\begin{mmlongtable}{Key terms for rapid revision}{tab:glossary3}%
{@{}P{0.26\textwidth}P{0.67\textwidth}@{}}
\rowcolor{ocre}\thd{Term} & \thd{One-line meaning}\\
\midrule
\endfirsthead
\rowcolor{ocre}\thd{Term} & \thd{One-line meaning}\\
\midrule
\endhead
\rlab{A/B testing} & Splitting traffic between two versions to find the
statistically better one\\
\rowcolor{tablealt}
\rlab{Acquisition channel} & The grouping describing where a session came from\\
\rlab{Auto-responder} & An automated sequence released to new subscribers\\
\rowcolor{tablealt}
\rlab{Bounce rate (web)} & Percentage of visitors who leave after viewing one
page without interaction\\
\rlab{Bounce rate (e-mail)} & Percentage of sent e-mails that were not
delivered\\
\rowcolor{tablealt}
\rlab{Buzz} & A combined popularity measure across social mentions\\
\rlab{CAN-SPAM} & US anti-spam law; fines up to \$16,000 per violation\\
\rowcolor{tablealt}
\rlab{Campaign analysis} & Evaluating e-mail performance on opens, clicks,
conversions, bounces and unsubscribes\\
\rlab{CRO} & Conversion rate optimisation --- systematically raising the
percentage who convert\\
\rowcolor{tablealt}
\rlab{CTR} & Percentage of recipients, or of opens, who clicked a link\\
\rlab{Dimension} & An attribute of the data --- city, device, source; forms the
rows of a report\\
\rowcolor{tablealt}
\rlab{Double opt-in} & Consent confirmed through a verification e-mail\\
\rlab{Drip campaign} & A scheduled nurture sequence delivered over time\\
\rowcolor{tablealt}
\rlab{E-mail marketing} & Using e-mail to deliver commercial and relationship
messages to a permission-based list\\
\rlab{Exit rate} & Percentage of sessions that ended on a given page\\
\rowcolor{tablealt}
\rlab{Funnel visualisation} & A report showing drop-off at each step of a defined
path\\
\rlab{Goal / conversion} & A defined, valuable user action being tracked\\
\rowcolor{tablealt}
\rlab{Google Analytics} & The platform that collects, measures and reports
website and app user data\\
\rlab{Google Tag Manager} & A container for deploying and versioning tracking
tags without code changes\\
\rowcolor{tablealt}
\rlab{List decay} & The natural annual loss of list quality, roughly 20--25\% per
year\\
\rlab{Metric} & A quantitative measurement --- sessions, revenue; forms the
columns of a report\\
\rowcolor{tablealt}
\rlab{Non-bounce rate} & $100\%$ minus bounced over total sent --- the delivery
rate\\
\rlab{Open rate} & Percentage of delivered e-mails opened\\
\rowcolor{tablealt}
\rlab{Opt-in} & Permission willingly given to receive e-mail\\
\rlab{Pageviews} & Total pages loaded; unique pageviews counts a page once per
session\\
\rowcolor{tablealt}
\rlab{Revenue per message} & Total revenue divided by the number of sent
e-mails\\
\rlab{Segment} & An isolated subset of users defined by shared attributes or
behaviour\\
\rowcolor{tablealt}
\rlab{Session} & A group of user interactions within a given time frame\\
\rlab{Share of voice} & Conversations about you against conversations about
competitors\\
\rowcolor{tablealt}
\rlab{Transactional e-mail} & A user-triggered informational e-mail --- receipt,
one-time password, shipping notice\\
\rlab{Trigger} & An automated e-mail fired by a specific user action\\
\rowcolor{tablealt}
\rlab{Unsubscribe rate} & Percentage of recipients who opted out\\
\rlab{Web analytics} & Collecting, measuring, analysing and reporting website
data to understand and optimise web usage\\
\end{mmlongtable}

\chapter{Framework Quick-Reference Sheet}

Everything in this volume that is a numbered list, in one place, for the night
before the examination.

\begin{mmlongtable}{Framework quick reference}{tab:quickref3}%
{@{}P{0.28\textwidth}P{0.65\textwidth}@{}}
\rowcolor{ocre}\thd{Framework} & \thd{The sequence}\\
\midrule
\endfirsthead
\rowcolor{ocre}\thd{Framework} & \thd{The sequence}\\
\midrule
\endhead
\rlab{Three metric families} & Volume (are we seen?), quality (are we relevant?),
action (are they engaging?)\\
\rowcolor{tablealt}
\rlab{Traditional against web analytics} & Data collection, speed, accuracy,
customer tracking, cost, measurement, sample, attribution, feedback loop\\
\rlab{Bounce against exit} & Bounce is a single-page no-interaction session; exit
is the last page of any session\\
\rowcolor{tablealt}
\rlab{Dimensions and metrics} & Attributes form the rows; numbers form the
columns\\
\rlab{Eight roles of web analytics} & Effectiveness, behaviour, conversion,
budget, content, technical, segmentation, forecasting\\
\rowcolor{tablealt}
\rlab{Seven benefits} & Conversion, acquisition cost, experience, problem
detection, content decisions, ROI proof, competitive advantage\\
\rlab{KPI alignment (five steps)} & Identify goals \ra{} map to KPIs \ra{} make
SMART \ra{} instrument analytics \ra{} optimise continuously\\
\rowcolor{tablealt}
\rlab{CRO process} & Find the weakest step \ra{} hypothesise \ra{} build the
variant \ra{} split traffic \ra{} run to significance \ra{} implement \ra{}
repeat\\
\rlab{Four GA features to name} & Traffic source analysis, bounce rate,
demographics, conversion and goal tracking\\
\rowcolor{tablealt}
\rlab{Nine acquisition channels} & Organic search, paid search, direct, referral,
organic social, paid social, e-mail, display, affiliates\\
\rlab{Seven report families} & Real-time, audience, acquisition, behaviour,
conversions, attribution, custom and explorations\\
\rowcolor{tablealt}
\rlab{Goal setup (eight steps)} & Define the conversion \ra{} choose the type
\ra{} implement \ra{} assign value \ra{} configure the funnel \ra{} verify \ra{}
track \ra{} act\\
\rlab{Tag Manager --- for and against} & Faster deployment, standard data layer,
consent, page weight, experimentation, rollback \emph{against} tag sprawl,
duplicate firing, performance, security\\
\rowcolor{tablealt}
\rlab{What against why} & Analytics locates the drop-off; heatmaps, recordings and
surveys explain it\\
\rlab{E-mail 1.0 against 2.0} & One-way, static, batch, untracked \emph{against}
triggered, segmented, automated, measured\\
\rowcolor{tablealt}
\rlab{Four e-mail advantages} & Low cost, global reach, easy tracking,
personalised communication\\
\rlab{Opt-in against spam} & Acquisition method, expectation setting, brand
impact, list quality\\
\rowcolor{tablealt}
\rlab{Three e-mail rules} & One topic, simple design complementing the message,
one clear call to action\\
\rlab{Nine e-mail types} & Revenue: promotional, new inventory, reorder.
Relationships: newsletter, welcome, educational, product advice. Intelligence:
testimonials, surveys\\
\rowcolor{tablealt}
\rlab{Transactional against promotional} & Trigger, purpose, consent, volume,
open rate (40--60\% against 15--25\%), examples, timing\\
\rlab{Seven-node e-mail plan} & Software \ra{} list \ra{} capture forms \ra{}
goals \ra{} auto-responders \ra{} triggers \ra{} metrics \ra{} back to software\\
\rowcolor{tablealt}
\rlab{Segmentation (eight steps)} & Collect data \ra{} choose criteria \ra{}
build dynamic segments \ra{} design per segment \ra{} personalise within \ra{}
automate triggers \ra{} test within \ra{} measure per segment\\
\rlab{Six-stage measurement funnel} & Sent \ra{} delivered \ra{} opened \ra{}
clicked \ra{} converted \ra{} revenue per message\\
\rowcolor{tablealt}
\rlab{Six e-mail challenges} & Low opens, spam filtering, unsubscribes, poor
design, no personalisation, inactive subscribers\\
\rlab{Offline tracking} & Point-of-sale e-mail capture, dedicated phone numbers,
e-mail-only entrance codes; promo codes, UTM tags, QR coupons\\
\rowcolor{tablealt}
\rlab{Analytics--e-mail integration} & Collect \ra{} segment \ra{} personalise
\ra{} track \ra{} retarget \ra{} optimise\\
\rlab{Four pillars} & The fuel (opt-in list), the payload (content arsenal), the
engine (triggers and tech), the dashboard (metrics)\\
\end{mmlongtable}


%%%%%%%%%%%%%%%%%%%%%%%%%%%%%%%%%%%%%%%%%%%%%%%%%%%%%%%%%%%%%%%%%%%%%%%%%%%%%%
\backmatter
\printindex

\end{document}
